\documentclass[10pt,a4paper]{amsart}
\usepackage[margin=19mm]{geometry}
\usepackage{amsmath,amssymb,mathtools}
\usepackage[scr=rsfs,cal=euler]{mathalfa}
\usepackage{xfrac}
\usepackage[unicode,hidelinks,bookmarks=false]{hyperref}
\newcommand{\ie}{i.e.\ }
\newcommand{\cf}{cf.\ }
\newcommand{\resp}{respectively\ }
\newcommand{\Subsection}{\subsection}
\providecommand{\nobreakdash}{\nobreak\hbox{-}}
\setlength{\emergencystretch}{2em}
\multlinegap0pt
\usepackage{mathrsfs}
\usepackage{url}
\providecommand{\doi}[1]{\textsc{doi:}\,\url{https://doi.org/#1}}
\usepackage{enumitem}
\newtheorem{theorem}{Theorem}[section]
\newtheorem{proposition}[theorem]{Proposition}
\newtheorem{lemma}[theorem]{Lemma}
\newtheorem{corollary}[theorem]{Corollary}
\newtheorem{conjecture}[theorem]{Conjecture}
\newtheorem{remark}[theorem]{Remark}
\newtheorem{example}[theorem]{Example}
\newtheorem{definition}[theorem]{Definition}
\newcommand{\Q}{\mathbb{Q}}
\newcommand{\Z}{\mathbb{Z}}
\newcommand{\F}{\mathbb{F}}
\newcommand{\C}{\mathbb{C}}
\newcommand{\R}{\mathbb{R}}
\newcommand{\GU}{\mathrm{GU}}
\newcommand{\GSp}{\mathrm{GSp}}
\newcommand{\Sp}{\mathrm{Sp}}
\newcommand{\SO}{\mathrm{SO}}
\newcommand{\Aut}{\mathrm{Aut}}
\newcommand{\Tr}{\mathrm{Tr}}
\newcommand{\Res}{\mathrm{Res}}
\newcommand{\Upr}{U_{\mathrm{pr}}}
\newcommand{\Va}{\mathrm{Va}}
\newcommand{\snew}{S^{\mathrm{new}}}
\newcommand{\eps}{\varepsilon}
\newcommand{\Onorm}{\mathcal{O}}
\newcommand{\one}{\mathbf{1}}
\DeclareMathOperator{\charp}{charpoly}
\DeclareMathOperator{\Jac}{Jac}
\DeclareMathOperator{\diag}{diag}
\DeclareMathOperator{\disc}{disc}
\DeclareMathOperator{\vol}{vol}
\DeclareMathOperator{\Orb}{O}
\usepackage{booktabs}
\providecommand{\kro}[2]{\left(\frac{#1}{#2}\right)}
\providecommand{\StatusKnown}{\textup{known}}
\providecommand{\StatusDerived}{\textup{derived}}
\providecommand{\StatusCertified}{\textup{certified}}
\providecommand{\Upr}{U_{\mathrm{pr}}}
\providecommand{\Va}{\mathrm{Va}}
\providecommand{\snew}{S^{\mathrm{new}}}
\providecommand{\charp}{\operatorname{charpoly}}
\usepackage{cleveref}
\begin{document}
\title[A structural trace identity and certified spectra for the Ri]{A structural trace identity and certified spectra for the Richelot-Brandt graph}
\author{Secondary 11F72, 11G10, 11R52, 14K02 ęywords Richelot isogeny graph, Brandt matrix, superspecial abelian surfaces, paramodular forms, Atkin--Lehner involution, trace formula ạte e̱gin document e̱gin abstract The degree-$2$ Brandt operator $B_2(2)$ on the principal genus of binary; Hung T. Dang}
\date{}
\maketitle
\noindent\emph{Source and licence.} A structural trace identity and certified spectra for the Richelot-Brandt graph. Version arXiv:2608.14145v2 (CC BY 4.0). This material is distributed under the Creative Commons Attribution 4.0 International licence, http://creativecommons.org/licenses/by/4.0/, as recorded in the arXiv OAI licence metadata for this exact version and shown on the abstract page https://arxiv.org/abs/2608.14145v2. Full rights notice: licenses/arxiv-2608.14145-CC-BY-4.0.txt.\par\smallskip
\noindent\textit{Editorial scope.} Counted unit: the original \texttt{theorem} environment \texttt{thm:occurrence} at source lines 1300--1343 together with its complete proof at lines 1345--1407; the delivered document keeps source lines 166--2650 so that every referenced label and every citation resolves inside it. Native editable source is the paper's own concatenated TeX; the AMS wrapper is editorial formatting only. No publisher class, upstream script or unrelated artwork is redistributed.
	\begin{theorem}[Structural trace identity]\label{thm:structural-trace-intro}
		For every prime $p\ge7$,
		\[
		\Tr\!\left(R(\pi)\mid M_{0,0}(\Upr(p))\right)
		=
		1+\dim\snew_4
		+\dim S_2^-\dim S_4^+
		-\dim S_2^+\dim S_4^-
		+d(p),
		\]
		where
		\[
		d(p)=\bigl(\dim S_3(K(p))-\dim S_4^-\bigr)-2\dim S_3^+(K(p)).
		\]
	\end{theorem}
	The full statement is given as \Cref{thm:structural-trace}  and its proof is in Section~\ref{sec:trace}. The proof starts from the closed principal- and non-principal-genus trace
	formulas of \cite{IbukiyamaKatsura,Ibukiyama2019quinary}. Their difference
	has a simple form: the coefficients of the terms involving
	$B_{2,\chi}$, $h(\sqrt{-2p})$, and $h(\sqrt{-3p})$ agree, while the
	remaining $h(\sqrt{-p})$ contribution is explicit. Combining this
	comparison with the non-principal-genus correspondence produces the stated
	principal-genus formula.
	
	The theorem fixes the signed total of the spectrum, but not its individual
	constituents. We formulate the following refinement of Ibukiyama's
	principal-genus multiplicity conjectures.
	
	\begin{conjecture}[Eigenvalue--sign refinement; verified for $11\le p\le149$]\label{conj:refined-intro}
		For every prime $p\ge11$, the characteristic polynomial of $B_2(2)$
		decomposes into an Eisenstein factor, Saito--Kurokawa factors,
		opposite-sign Yoshida factors, an exact square
		$\mathrm{Va}_p(x)^2$, and a general-type factor $N_p(x)$. The involution
		$R(\pi)$ has the predicted signs on the Eisenstein, Saito--Kurokawa,
		Yoshida, and general-type components, and each type-$\mathrm{Va}$ pair
		splits into one $+1$ and one $-1$ eigenline.
	\end{conjecture}
	The full statement is given as Conjecture~\ref{conj:eigenvalue-sign} in Section~\ref{sec:statements}. The eigenvalue formulas entering this conjecture are consequences of the transfer to $\GSp_4$ and the spin normalization at $2$: an occurring
	Eisenstein, Saito--Kurokawa, Yoshida, or general-type parameter contributes
	respectively an eigenvalue
	\[
	15,\qquad a_2(g)+6,\qquad a_2(g)+2a_2(f),\qquad a_2(F).
	\]
	The accompanying decomposition of the total dimension is equivalent to
	Ibukiyama's unconditional dimension identity; no new dimension theorem is
	claimed here.
	
	The type-Va clause of the conjecture is a theorem for every prime.
	
	\begin{theorem}[Occurrence and splitting of the type-Va pairs]\label{thm:occurrence-intro}
		For every prime $p\ge11$ and every general-type parameter whose local
		component at $p$ is of type $\Va$, both members of the local
		$L$-packet at $p$ occur in $M_{0,0}(\Upr(p))$, each with automorphic
		multiplicity one; the resulting isotypic block has even dimension,
		$R(\pi)$-trace zero, and $+1$ and $-1$ eigenspaces of equal
		dimension. In particular the type-$\Va$ contribution to the
		characteristic polynomial of $B_2(2)$ is an exact square split
		evenly by $R(\pi)$.
	\end{theorem}
	The full statement is Theorem~\ref{thm:occurrence} in
	Section~\ref{sec:occurrence}. The input is the global lifting theorem of
	R\"osner and Weissauer for inner forms of $\GSp_4$ anisotropic at the
	archimedean place \cite[Thm.~11.4]{RoesnerWeissauer2021}, which applies to
	$\GU_2(B)$ at trivial weight; the $\pm$ separation then follows from the
	local Langlands correspondence for $\GU_2(B_p)$ as in
	Remark~\ref{rem:localpacket}. What remains conjectural in
	Conjecture~\ref{conj:refined-intro} is the occurrence pattern of the
	Saito--Kurokawa and Yoshida rows and the resulting identification of the
	block degrees.
	
	
	\begin{theorem}[Certified spectra through $p=149$]\label{thm:certified-intro}
		For every prime $11\le p\le149$, the exact-arithmetic certificates in the
		archived Richelot--Brandt artifact verify the factorization and the
		$R(\pi)$-sign decomposition predicted by
		Conjecture~\ref{conj:refined-intro}.
	\end{theorem}
	
	The full statement is given as Theorem~\ref{thm:certified-spectra}, and its proof is in Section~\ref{sec:certificates}. The proof of this finite statement has two parts. The geometric pipeline produces the Richelot matrix, its weights, and the $R(\pi)$-permutation.
	The elliptic pipelines independently supply the weight-$2$ and weight-$4$
	block data. Exact polynomial division and isotypic projector traces then
	verify the predicted factorization and signs. The trace identity above is
	an additional global consistency check, not a replacement for these
	factor-by-factor certificates.
	
	\subsection{What is new}\label{ss:what-is-new}
	
	The contribution has four distinct components, each with its own logical
	status.
	
	\begin{enumerate}[label=\textup{(\roman*)},leftmargin=2.5em]
		\item The trace comparison leading to
		Theorem~\ref{thm:structural-trace-intro}, and the structural
		principal-genus form of the resulting trace identity.
		\item The eigenvalue--sign formulation of
		Conjecture~\ref{conj:refined-intro}, especially the use of $R(\pi)$ to
		separate the type-$\mathrm{Va}$ pairs that cannot be separated by
		dimension data or by Hecke operators away from $p$.
		\item The certified finite verification through $p=149$, which constructs
		the graph-side data geometrically and checks the predicted decomposition in
		exact arithmetic.
		\item The occurrence theorem for the type-$\Va$ block
		(Section~\ref{sec:occurrence}): by the global lifting theorem of
		R\"osner and Weissauer, both members of every type-$\Va$ pair occur
		with multiplicity one, so the square and sign assertions of
		Conjecture~\ref{conj:eigenvalue-sign}\textup{(iii)} are theorems for
		every prime.
	\end{enumerate}
	The third component is a methodological one as much as a computational one.
	Every step of the verification is carried out in exact arithmetic:
	characteristic polynomials are computed over $\Z$ and divided exactly by the
	block polynomials, and the isotypic signs are decoded from projector traces
	modulo a single prime whose sufficiency is proved rather than assumed, the
	bound $2h_2(p)\le2866<2^{61}-1$ of Lemma~\ref{lem:onemod} being what makes
	one modulus enough. The archive of \S\ref{ss:reproducibility} is arranged so
	that the certificates can be replayed from the frozen records alone, by a
	verifier that shares no code with the engines that produced them.
	The local language suggested by the trace comparison must be interpreted
	carefully. The proof is a global comparison of closed trace formulas. It
	does not provide a direct termwise proof of equality between local twisted
	orbital integrals at the two non-conjugate maximal parahorics. We regard a
	direct local proof as a separate problem.
	
	Of these four components, the first and the fourth are unconditional for all $p$; the second is a conjecture whose remaining content is the Saito--Kurokawa and Yoshida rows and the identification of the block degrees; the third certifies the full conjecture in a finite, explicitly bounded range and does not constitute a proof beyond it.
	
	\subsection{Relation to prior work}\label{ss:relation-prior}
	
	The transfer and Arthur-type framework used here rests on published work on
	the compact inner form and its transfer to $\GSp_4$
	\cite{RoesnerWeissauer2021,vanHoften,DPRT,Schmidt2018}. The parahoric data
	that identify the type-$\mathrm{Va}$ contribution come from
	\cite[Table~A.15]{RobertsSchmidt}. Ibukiyama's dimension identity supplies
	the unconditional count underlying the degree bookkeeping
	\cite[Thm.~3.1]{Ibukiyama2018conj}; the multiplicity assertion for the
	general-type rows is proved in Section~\ref{sec:occurrence} via
	\cite[Thm.~11.4]{RoesnerWeissauer2021}, while the Saito--Kurokawa and
	Yoshida rows remain conjectural outside the certified range.
	
	The paper is complementary to work on spectral gaps and expansion in
	higher-dimensional isogeny graphs. Those results control the size or
	distribution of the spectrum. Here the aim is instead to identify the
	source of individual factors of a single weighted Brandt spectrum. The
	resulting factorization makes visible the Saito--Kurokawa, Yoshida,
	type-$\mathrm{Va}$, and general-type contributions separately.
	
	\subsection{Logical guide}\label{ss:logical-guide}
	The four levels above recur throughout the paper under their formal names; the following table fixes the correspondence.
	
	\begin{center}
		\begin{tabular}{lll}
			\toprule
			Statement & Scope & Status \\
			\midrule
			Structural trace identity & $p\ge7$ & unconditional theorem \\
			Dimension bookkeeping & $p\ge11$ & reformulation of a known theorem \\
			Transfer and eigenvalue laws & $p\ge11$ & derived from published results \\
			Eigenvalue--sign refinement & $p\ge11$ & conjecture \\
			Certified spectra & $11\le p\le149$ & computer-assisted theorem \\
			Fricke-sign bias $d(p)\ge0$ & every prime & unconditional theorem \\
			\bottomrule
		\end{tabular}
	\end{center}
	The same list, with the internal numbering used in the body of the paper,
	is collected as Table~\ref{tab:logical-status} in
	Section~\ref{sec:statements}.
	
	The final theorem relies on frozen exact data and a verifier. External
	comparisons with paramodular data, including the quinary database, are
	reported as corroboration and are not inputs to its proof.
	
	\smallskip
	\noindent\emph{Organization.}
	Section~\ref{sec:prelim} fixes the arithmetic and geometric settings and
	all normalizations. The transfer framework and the conjectural
	eigenvalue--sign refinement are developed next. The trace comparison and
	structural trace identity are then proved. The certified computation is
	presented separately, including its certificate format and reproducibility
	protocol. The remaining sections derive consequences for type numbers,
	Ihara pairs, non-lift factors, and the Fricke-sign bias.
	
	\section{Statement of results and logical status}
	\label{sec:statements}
	
	This section states the principal results and records their logical status
	before introducing the notation used in the proofs. The paper separates
	four levels of assertion: an unconditional trace theorem, a conjectural
	eigenvalue--sign refinement, a finite certified computation, and a
	transfer framework derived from published results.
	
	For $k\in\{2,4\}$, let $\snew_k=\snew_k(\Gamma_0(p))$ denote the space of
	elliptic newforms of weight $k$ and prime level $p$, and write
	$S_k^{\pm}$ for its Atkin--Lehner eigenspaces. We use the following
	combinations throughout:
	\begin{align}
		\mathrm{cross}(p)
		&:=\dim S_2^-\dim S_4^+ + \dim S_2^+\dim S_4^-,
		\label{eq:cross-def}\\
		\delta(p)
		&:=\dim S_3(K(p))-\dim S_4^-,
		\label{eq:delta-def}\\
		d(p)
		&:=\delta(p)-2\dim S_3^+(K(p)),
		\label{eq:d-def}\\
		L_0(p)
		&:=1+\dim\snew_4+\dim S_2^-\dim S_4^+-\dim S_2^+\dim S_4^-.
		\label{eq:L0-def}
	\end{align}
	Here $\delta(p)$ counts the weight-$3$ paramodular non-lifts, while $d(p)$
	counts them with Fricke signs. The lift part $L_0(p)$ is the portion of the
	trace accounted for by Eisenstein, Saito--Kurokawa, and Yoshida classes.
	
	\subsection{The structural trace identity}
	\begin{theorem}[Structural trace identity; full form of \Cref{thm:structural-trace-intro}]\label{thm:structural-trace}
		For every prime $p\ge7$,
		\begin{equation}\label{eq:structural-trace}
			\Tr\!\left(R(\pi)\mid M_{0,0}(\Upr(p))\right)=L_0(p)+d(p).
		\end{equation}
		Equivalently,
		\[
		\Tr\!\left(R(\pi)\mid M_{0,0}(\Upr(p))\right)
		=
		1+\dim\snew_4
		+\dim S_2^-\dim S_4^+
		-\dim S_2^+\dim S_4^-
		+d(p).
		\]
	\end{theorem}
	
	The proof is a termwise comparison of the closed principal- and
	non-principal-genus trace evaluations of
	\cite{IbukiyamaKatsura,Ibukiyama2019quinary}, combined with the
	non-principal-genus correspondence of \cite{DPRT,Ibukiyama-dim}. The
	comparison is global. It does not supply a direct local proof of equality
	between twisted orbital integrals at the two maximal parahorics.
	
	The following comparison is the key intermediate identity. Let
	$R_{\mathrm{npg}}(\pi)$ denote the corresponding involution on the
	non-principal genus, let $e_3=\bigl(\frac p3\bigr)$, and let $\nu(p)$ be
	the class-number quantity defined in \S\ref{sec:trace}. Then
	\begin{equation}\label{eq:genus-comparison}
		\Tr R(\pi)-\Tr R_{\mathrm{npg}}(\pi)
		=\nu(p)\frac{p-e_3}{12}.
	\end{equation}
	The coefficients of $B_{2,\chi}$, $h(\sqrt{-2p})$, and $h(\sqrt{-3p})$
	agree in the two evaluations; only the contribution represented by
	$h(\sqrt{-p})$ remains.
	
	\subsection{The eigenvalue--sign refinement}
	
	The trace theorem fixes the signed total of the anticipated spectrum, but
	not its individual constituents. The next statement is conjectural in
	general.
	
	\begin{conjecture}[C1: eigenvalue--sign refinement; verified for $11\le p\le149$]
		\label{conj:eigenvalue-sign} \label{conj:refined}
		For every prime $p\ge11$,
		\begin{equation}\label{eq:C1-factorization}
			\begin{split}
				\charp B_2(2)(x)
				={}&(x-15)
				\prod_{g\subset\snew_4}m_g(x-6)
				\prod_{\substack{(f,g):\\ \eps_p(f)\eps_p(g)=-1}}
				\Res_y\bigl(t_f(y),m_g(x-y)\bigr)\\
				&\qquad\qquad\cdot \Va_p(x)^2N_p(x).
			\end{split}
		\end{equation}
		Here $m_g$ is the minimal polynomial of $a_2(g)$ and $t_f$ is the minimal
		polynomial of $2a_2(f)$. Moreover:
		\begin{enumerate}[label=\textup{(\roman*)},leftmargin=2.5em]
			\item the Saito--Kurokawa product runs over all weight-$4$ newform orbits;
			\item the Yoshida product runs over the opposite-sign pairs;
			\item the type-$\Va$ block is an exact square, and each such pair splits
			under $R(\pi)$ into one $+1$ and one $-1$ eigenline;
			\item $N_p$ is monic of degree $\delta(p)$, with roots the spin
			$T(2)$-eigenvalues of the weight-$3$ paramodular non-lifts;
			\item the $R(\pi)$-signs are $+1$ on the Eisenstein and
			Saito--Kurokawa components, $-\eps_p(f)$ on the Yoshida component attached
			to $(f,g)$, and $-\eps_p(F)$ on the general-type component attached to a
			non-lift $F$.
		\end{enumerate}
	\end{conjecture}
	
	The conjecture makes two logically distinct assertions: an occurrence and
	multiplicity assertion, namely that the residual factor is an exact square
	of the stated degree, and a sign assertion about $R(\pi)$ on each block.
	For the type-Va block both are proved for every prime in
	Section~\ref{sec:occurrence} (Theorem~\ref{thm:occurrence} and
	Corollary~\ref{cor:conj-iii}), up to the identification of the degree,
	which rests on the remaining rows of Table~\ref{tab:twocol}.
	Theorem~\ref{thm:certified-spectra} establishes the full statement in the
	certified range.
	
	This refinement is strictly stronger than a dimension statement. Unlike
	the Saito--Kurokawa and Yoshida blocks, $\Va_p$ is not given by elliptic
	data: it is the residual left by the bookkeeping identity, and its roots
	are eigenvalues of $B_2(2)$ that no weight-$2$ or weight-$4$ newform
	predicts. The two
	members of a type-$\Va$ pair have the same eigenvalue for every $T(n)$
	with $p\nmid n$, so neither dimensions nor the away-from-$p$ Hecke
	algebra separates them. Their separation is visible only through the local
	involution $R(\pi)$.
	
	\subsection{The certified finite verification}
	
	\begin{theorem}[Certified spectra through $p=149$]
		\label{thm:certified-spectra}
		For every prime $11\le p\le149$, the archived exact-arithmetic data and the
		accompanying verifier satisfy the certificate hypotheses of \S\ref{sec:certificates}.
		Consequently, Conjecture~\ref{conj:eigenvalue-sign} holds for each of these thirty-one
		primes.
	\end{theorem}
	
	The certificates check the weighted graph relations for the Richelot
	matrix, the involution, exact division by the Eisenstein,
	Saito--Kurokawa, and Yoshida blocks, the exact square in the residual
	factor, and the signed isotypic multiplicities determined by projector
	traces. The statement is finite and computer-assisted; it does not prove
	Conjecture~\ref{conj:eigenvalue-sign} beyond the certified range.
	
	Two cases display the content concretely. At $p=19$, the factor
	$(x+2)^2$ is the first type-$\Va$ block in the certified range and is
	split by $R(\pi)$. At $p=61$,
	\[
	N_{61}(x)=x+7,
	\]
	so the graph realizes the simple general-type eigenvalue $-7$.
	
	\subsection{Transfer and eigenvalue framework}
	
	The following theorem explains the eigenvalue laws used in the
	conjectural factorization.
	
	\begin{theorem}[Transfer and eigenvalue laws]\label{thm:transfer-framework}
		Let $p\ge11$. Every simultaneous eigensystem of $\{B_2(2),R(\pi)\}$ on
		$M_{0,0}(\Upr(p))$ transfers to a split-side Arthur parameter of type
		\textup{(Eis)}, \textup{(P)}, \textup{(Y)}, or \textup{(G)}, with the
		same unramified Hecke data at every $q\nmid p$. In the normalization
		\[
		Q_2(X)=1-\lambda_2X+\cdots+2^6X^4,
		\]
		the corresponding $B_2(2)$-eigenvalue is respectively
		\[
		15,\qquad a_2(g)+6,\qquad a_2(g)+2a_2(f),\qquad a_2(F).
		\]
		The spectrum is real, $15$ is a simple eigenvalue of $R(\pi)$-sign $+1$,
		and every other eigenvalue lies in $(-15,15)$.
	\end{theorem}
	
	\noindent Theorem~\ref{thm:transfer-framework} is proved in
	\S\ref{sec:prelim}--\S\ref{sec:transfer}: the transfer and the eigenvalue
	formulas are Proposition~\ref{prop:arthur}, and the spectral clauses are
	Proposition~\ref{prop:frame} together with Lemma~\ref{lem:eis}.
	
	The associated block-degree bookkeeping is
	\begin{equation}\label{eq:dimension-bookkeeping}
		h_2(p)=1+\dim\snew_4+\mathrm{cross}(p)+2\deg\Va_p+\delta(p).
	\end{equation}
	Under the change of variables in \S\ref{sec:transfer}, this is equivalent
	to Ibukiyama's dimension theorem \cite[Thm.~3.1]{Ibukiyama2018conj}. It is
	therefore a reformulation of known dimension information, not a new
	class-number theorem.
	
	\subsection{Logical status}
	Table~\ref{tab:logical-status} summarizes the logical status of the
	principal assertions.
	
	\begin{table}[ht]
		\centering
		\begin{tabular}{@{}p{0.34\linewidth}p{0.17\linewidth}p{0.35\linewidth}@{}}
			\toprule
			Statement & Range & Status \\
			\midrule
			Theorem~\ref{thm:structural-trace} & $p\ge7$ & Unconditional theorem \\
			Equation~\eqref{eq:dimension-bookkeeping} & $p\ge11$ & Reformulation of a known dimension theorem \\
			Theorem~\ref{thm:transfer-framework} & $p\ge11$ & Derived from published transfer and classification results \\
			Theorem~\ref{thm:occurrence} & $p\ge11$ & Unconditional theorem, via \cite[Thm.~11.4]{RoesnerWeissauer2021} \\
			Conjecture~\ref{conj:eigenvalue-sign} & $p\ge11$ & Type-Va clause proved (Theorem~\ref{thm:occurrence}); remaining rows conjectural \\
			Theorem~\ref{thm:certified-spectra} & $11\le p\le149$ & Finite computer-assisted theorem with archived certificates \\
			Theorem~\ref{thm:bias} & every prime & Unconditional theorem, via the exact formula for $d(p)$ \\
			\bottomrule
		\end{tabular}
		\caption{Logical status of the principal assertions.}
		\label{tab:logical-status}
	\end{table}
	
	Theorem~\ref{thm:structural-trace} is compatible with
	Conjecture~\ref{conj:eigenvalue-sign} at the level of signed totals, but
	it does not determine the individual multiplicities asserted by the
	conjecture. Theorem~\ref{thm:certified-spectra} verifies those
	multiplicities only in the finite range. External comparisons with
	paramodular databases are reported separately as corroboration and are not
	inputs to the proofs. The exact closed formula for the defect $d(p)$ and
	its positivity (Theorem~\ref{thm:bias}) are proved in the final section.
	
	\section{Richelot--Brandt setting and normalizations}
	\label{sec:prelim}
	
	This section fixes the arithmetic, geometric, and automorphic settings and
	all conventions used later. Unless stated otherwise, vector spaces are over
	$\C$ and $p$ is a prime.
	
	\subsection{The principal genus and its operators}
	\label{ss:principal-genus}
	
	Let $B=B_{p,\infty}$ be the definite quaternion algebra over $\Q$ ramified
	exactly at $p$ and $\infty$, with canonical involution $x\mapsto\bar x$,
	reduced trace $\mathrm{tr}\,x=x+\bar x$, and reduced norm
	$\mathrm{n}\,x=x\bar x$. Fix a maximal order $\Onorm\subset B$. By
	Deuring's correspondence, the supersingular elliptic curves over
	$\overline{\F}_p$ correspond to the left $\Onorm$-ideal classes, of number
	$h=h(p)$, and their masses satisfy Eichler's relation
	\begin{equation}\label{eq:eichler-mass}
		\sum_{E}\frac{1}{\#\Aut E}=\frac{p-1}{24}.
	\end{equation}
	
	On $B^2$ the standard positive definite quaternion Hermitian form \cite{Shimura1963}
	$H(x,y)=\bar x_1y_1+\bar x_2y_2$ has unitary similitude group
	\[
	G=\GU_2(B),
	\]
	an inner form of $\GSp_4$ over $\Q$ \cite{Ibukiyama2007compact}. The left
	$\Onorm$-lattices in $B^2$ fall into genera and classes. The principal
	genus of the standard lattice $\Onorm^2$ has class number $h_2=h_2(p)$,
	given in closed form by \cite{HashimotoIbukiyama}. Locally at $p$, the two
	genera of maximal type correspond to two conjugacy classes of maximal open
	compact subgroups of $G(\Q_p)$, of different types and not conjugate; we
	denote them $U_1$ (principal) and $U_2$ (non-principal).
	
	Let $L_1,\ldots,L_{h_2}$ be representatives of the principal-genus classes.
	We write
	\[
	M_{0,0}(\Upr(p))=\{f:\{L_1,\ldots,L_{h_2}\}\to\C\}
	\]
	with the mass inner product
	\begin{equation}\label{eq:mass-inner}
		\langle f,g\rangle
		=\sum_{i=1}^{h_2}\frac{f(L_i)\,\overline{g(L_i)}}{e_i},
		\qquad e_i=\#\Aut(L_i).
	\end{equation}
	
	The Brandt matrices $B_n(2)$ of \cite{Hashimoto1980} act on this space and
	realize the Hecke operators $T(n)$; we abbreviate $B_2(2)$ for the
	degree-$2$ operator. The unique maximal two-sided ideal $\pi$ of $\Onorm$
	of reduced norm $p$ satisfies $\pi^2=p\,\Onorm$, so the induced operator
	$R(\pi)$ is an involution,
	\[
	R(\pi)^2=\mathrm{id},
	\]
	commuting with every $T(n)$. It is the Atkin--Lehner involution of the
	principal genus. When the non-principal genus is used, the corresponding
	operator is written $R_{\mathrm{npg}}(\pi)$.
	
	\subsection{The Richelot graph}
	\label{ss:richelot-graph}
	
	A principally polarized abelian surface over $\overline{\F}_p$ is either
	the Jacobian of a smooth genus-$2$ curve with its theta polarization, or a
	polarized product of two elliptic curves. It is superspecial if, after
	forgetting the polarization, it is isomorphic to a product of supersingular
	elliptic curves. By \cite[\S1]{IKO1986}, the superspecial principally
	polarized abelian surfaces over $\overline{\F}_p$, with their
	automorphisms, are exactly the classes of the principal genus; hence the
	vertex set of the graph below is identified with the index set of
	\eqref{eq:mass-inner}.
	
	A Richelot $(2,2)$-isogeny is an isogeny of principally polarized abelian
	surfaces whose kernel is a maximal Weil-isotropic subgroup of the
	$2$-torsion. Each source has exactly
	\begin{equation}\label{eq:fifteen}
		15=(2^2+1)(2+1)=\frac{6!}{2^3\,3!}
	\end{equation}
	such kernels: the number of maximal isotropic subspaces of a
	$4$-dimensional symplectic $\F_2$-space, equivalently the number of
	partitions of six Weierstrass points into three unordered pairs. These
	kernels realize $B_2(2)$ \cite{KatsuraTakashima}.
	
	Concretely, for a curve $C:y^2=f(x)$ with $f=c\,q_1q_2q_3$ a factorization
	into three quadratics over a splitting field, the codomain is controlled by
	the Richelot determinant
	\[
	\Delta_R=\det(q_1,q_2,q_3),
	\]
	the determinant of the coefficient matrix. If $\Delta_R\ne0$, the target is
	the Jacobian of
	\[
	\Delta_R\,y^2=[q_2,q_3]\,[q_3,q_1]\,[q_1,q_2],
	\qquad [q_i,q_j]=q_i'q_j-q_iq_j';
	\]
	if $\Delta_R=0$, the target is a product of elliptic curves
	\cite[\S8]{JordanZaytman}. The fifteen unordered pair-partitions of the
	roots of $f$ enumerate the kernels.
	
	\subsection{The dictionary and the Mestre symmetry}
	\label{ss:dictionary}
	
	We use the out-edge convention throughout:
	\begin{equation}\label{eq:out-edge}
		M_{ij}=\#\{\text{kernels at vertex $i$ with target $j$}\},
	\end{equation}
	so every row of $M$ sums to $15$. With $e_i=\#\Aut(A_i,\lambda_i)$ the
	vertex weights, the Mestre symmetry reads
	\begin{equation}\label{eq:mestre}
		e_jM_{ij}=e_iM_{ji}\qquad\text{for all }i,j,
	\end{equation}
	which is \cite[Thm.~18(c)]{JordanZaytman}.
	
	\begin{lemma}\label{lem:mestre}
		Let $M$ be a real $h\times h$ matrix, let $e_1,\ldots,e_h>0$ satisfy
		\eqref{eq:mestre}, and set $D=\diag(e_i)$. Then
		\begin{enumerate}[label=\textup{(\alph*)},leftmargin=2.2em]
			\item $S:=D^{-1/2}MD^{1/2}$ is symmetric;
			\item $M^{T}=D^{-1}MD$;
			\item $\charp(M)=\charp(S)$, so $M$ is diagonalizable with real spectrum;
			\item $M$ is self-adjoint for
			$\langle u,v\rangle=\sum_i u_iv_i/e_i$.
		\end{enumerate}
		Conversely, each of \textup{(a)} and \textup{(d)} is equivalent to
		\eqref{eq:mestre}. Moreover, the symmetrized support
		$\{(i,j):S_{ij}>0\}$ is exactly the undirected edge set of $G_p$:
		$S_{ij}>0\iff M_{ij}>0\iff M_{ji}>0$.
	\end{lemma}
	
	\begin{proof}
		$S_{ij}=e_i^{-1/2}M_{ij}e_j^{1/2}$, so $S_{ij}=S_{ji}\iff
		e_jM_{ij}=e_iM_{ji}$, giving (a) and its converse. For (b),
		$(D^{-1}MD)_{ij}=e_i^{-1}M_{ij}e_j=M_{ji}$. For (c), $S$ is conjugate to
		$M$ and symmetric, hence real-diagonalizable. For (d), for real $M$,
		\[
		\langle Mu,v\rangle-\langle u,Mv\rangle
		=\sum_{i,j}\bigl(M_{ij}/e_i-M_{ji}/e_j\bigr)u_jv_i,
		\]
		which vanishes identically iff \eqref{eq:mestre} holds. For the support
		claim, $S_{ij}=e_i^{-1/2}e_j^{1/2}M_{ij}$ with $e_i>0$, so
		$S_{ij}>0\iff M_{ij}>0$; and $M_{ij}>0\iff M_{ji}>0$ by \eqref{eq:mestre}.
	\end{proof}
	
	The last clause of Lemma~\ref{lem:mestre} is what makes the
	non-bipartiteness argument below valid: the sign pattern of a hypothetical
	$(-15)$-eigenvector is a $2$-colouring of the undirected graph whose edges
	are precisely the positive entries of $S$.
	
	The weights and the decomposable locus are pinned by the following lemma.
	
	\begin{lemma}\label{lem:vertices}
		Let $w_a=\#\Aut(E_a)/2\in\{1,2,3\}$ for the supersingular curves $E_a$
		over $\overline{\F}_p$, $a=1,\ldots,h$.
		\begin{enumerate}[label=\textup{(\alph*)},leftmargin=2.2em]
			\item The product vertices of $G_p$ are the unordered pairs
			$\{E_a,E_b\}$; their number is $h(h+1)/2$.
			\item Their weights are $e_{\{a,b\}}=\#\Aut(E_a\times E_b)$, equal to
			$8w_a^2$ if $a=b$ and $4w_aw_b$ if $a\ne b$.
			\item $\sum_a 1/w_a=(p-1)/12$, equivalently the Eichler mass
			$\sum_a 1/\#\Aut(E_a)=(p-1)/24$ of \eqref{eq:eichler-mass}.
			\item $G_p$ is connected and the Perron eigenvalue $15$ of $B_2(2)$ is
			simple.
		\end{enumerate}
	\end{lemma}
	
	\begin{proof}
		(a), (b) A polarized product $E_a\times E_b$ is superspecial, and two
		products are isomorphic as principally polarized abelian surfaces if and
		only if the unordered pairs of $j$-invariants agree \cite[\S3]{IKO1986};
		the automorphism group of the polarized product is
		$(\Aut E_a\times\Aut E_b)\rtimes C_2$ when $a=b$, of order
		$2(2w_a)^2=8w_a^2$, and $\Aut E_a\times\Aut E_b$ when $a\ne b$, of order
		$(2w_a)(2w_b)=4w_aw_b$, giving the stated weights. (c) Since
		$w_a=\#\Aut E_a/2$, $\sum_a 1/w_a=2\sum_a 1/\#\Aut E_a=(p-1)/12$ by
		\eqref{eq:eichler-mass}. (d) Connectivity is the irreducibility of the
		Brandt graph \cite[Thm.~44]{JordanZaytman}; a connected non-negative
		symmetrizable matrix with constant row sum $15$ has $15$ as a simple
		Perron eigenvalue by Perron--Frobenius applied to the symmetrization $S$
		of Lemma~\ref{lem:mestre}.
	\end{proof}
	
	\subsection{The spectral frame}
	\label{ss:spectral-frame}
	
	\begin{proposition}\label{prop:frame}
		For every prime $p\ge11$ the multigraph $G_p$ is connected and not
		bipartite; consequently $x-15$ divides $\charp B_2(2)$ with multiplicity
		one, and every other eigenvalue lies in the open interval $(-15,15)$. In
		particular any integer eigenvalue $\lambda\ne15$ satisfies
		$|\lambda|\le14$.
	\end{proposition}
	
	\begin{proof}
		Connectedness and non-bipartiteness of $G_p$ are
		\cite[Thm.~44]{JordanZaytman}. Let $S$ be the symmetrization of
		Lemma~\ref{lem:mestre}: $S$ is symmetric, nonnegative, and irreducible by
		connectedness, with row sums $15$, so its spectral radius is the simple
		Perron eigenvalue $15$. If $Sv=-15v$ with $v\ne0$, then entrywise
		$S|v|\ge|Sv|=15|v|$. Let $w>0$ be the Perron eigenvector, $Sw=15w$.
		Pairing and using $S=S^{T}$,
		\[
		15\,\langle w,|v|\rangle=\langle Sw,|v|\rangle=\langle w,S|v|\rangle
		\ge15\,\langle w,|v|\rangle,
		\]
		so $\langle w,S|v|-15|v|\rangle=0$; as $w>0$ and $S|v|-15|v|\ge0$
		entrywise, this forces $S|v|=15|v|$. Equality in $|(Sv)_i|=(S|v|)_i$ on
		each row then forces $v_j$ to have constant sign along the support of row
		$i$, i.e.\ $v_i$ and $v_j$ have opposite signs whenever $S_{ij}>0$. By the
		last clause of Lemma~\ref{lem:mestre} the set $\{(i,j):S_{ij}>0\}$ is the
		undirected edge set of $G_p$, so $\operatorname{sgn}v$ is a proper
		$2$-colouring of $G_p$, contradicting non-bipartiteness. Hence $-15$ is
		not an eigenvalue, and since the spectrum is real every non-Perron
		eigenvalue lies in $(-15,15)$.
	\end{proof}
	
	\subsection{Elliptic-modular data on the split side}
	\label{ss:split}
	
	For $k\in\{2,4\}$ let $\snew_k$ denote the space of weight-$k$ level-$p$
	elliptic newforms, and for $f\in\snew_k$ let $\eps_p(f)\in\{\pm1\}$ be its
	Atkin--Lehner eigenvalue at $p$, with $S_k^{\pm}$ the corresponding
	eigenspaces. At exact prime level the $p$-th Hecke eigenvalue is determined
	by the sign,
	\begin{equation}\label{eq:ap-sign}
		a_p(f)=-\eps_p(f)\,p^{k/2-1},
	\end{equation}
	by Atkin--Lehner theory \cite[Thm.~3]{AtkinLehner}.
	
	These spaces enter the catalogue through the Yoshida lifts: a Galois orbit
	$g\subset\snew_4$ and a form $f\in\snew_2$ contribute a paramodular
	Yoshida packet whose spin $T(2)$-eigenvalue is $a_2(g)+2a_2(f)$.
	Accordingly we set $m_g\in\Z[x]$ to be the minimal polynomial of $a_2(g)$
	and $t_f\in\Z[x]$ that of $2a_2(f)$, so that the monic resultant
	$\Res_y\bigl(t_f(y),m_g(x-y)\bigr)$ has root set exactly
	$\{a_2(g)+2a_2(f)\}$.
	
	\begin{lemma}\label{lem:resultant}
		Let $t,m\in\Z[x]$ be monic of degrees $a,b$ with complex roots
		$\alpha_i,\beta_j$. With
		$\Res_y(A,B)=\mathrm{lc}(A)^{\deg B}\prod_{A(\alpha)=0}B(\alpha)$,
		\[
		\Res_y\bigl(t(y),\,m(x-y)\bigr)
		=\prod_{i,j}\bigl(x-(\alpha_i+\beta_j)\bigr)\in\Z[x]
		\]
		is monic of degree $ab$ with root set $\{\alpha_i+\beta_j\}$, independent
		of the ordering convention up to the sign $(-1)^{ab}$ which the monic
		normalization removes. For monic $P,Q\in\Z[x]$, $\gcd_{\Q[x]}(P,Q)=1$ if
		and only if $P,Q$ share no complex root.
	\end{lemma}
	
	\begin{proof}
		Monicity of $t$ gives $\Res_y(t,m(x-\cdot))=\prod_i m(x-\alpha_i)$; each
		factor is monic of degree $b$ with roots $\alpha_i+\beta_j$. Coefficients
		are symmetric functions of the $\alpha_i$ with integer values, hence in
		$\Z$. The gcd statement is the nullstellensatz in one variable.
	\end{proof}
	
	\subsection{Spin normalization at $2$}
	\label{ss:spin-normalization}
	
	The eigenvalue laws of Theorem~\ref{thm:transfer-framework} use the
	normalization of the spin parameter at $2$ recorded in
	Table~\ref{tab:spin-normalization}.
	
	\begin{table}[ht]
		\centering
		\begin{tabular}{@{}lll@{}}
			\toprule
			Arthur type & Spin roots at $2$ & $B_2(2)$-eigenvalue \\
			\midrule
			Eisenstein & $\{1,2,4,8\}$ & $15$ \\
			Saito--Kurokawa, $g\subset\snew_4$ & $\{\alpha_g,\beta_g,2,4\}$ & $a_2(g)+6$ \\
			Yoshida, pair $(f,g)$ & $\{2\alpha_f,2\beta_f,\alpha_g,\beta_g\}$ & $a_2(g)+2a_2(f)$ \\
			General type & irreducible spin & $a_2(F)$ \\
			\bottomrule
		\end{tabular}
		\caption{Spin normalization at $2$ and the resulting eigenvalue laws.}
		\label{tab:spin-normalization}
	\end{table}
	
	The multiplier on each $\mathrm{GL}_2$-slot is forced by the
	central-character constraint: the product of the four spin roots equals
	$2^{4k-6}=64$ at $k=3$. This gives $16m^2=64$, i.e.\ $m=2$, for the
	weight-$2$ slot of type (Y), and the two Saito--Kurokawa slots of type (P)
	carry the Satake multipliers $2^{k-2}$ and $2^{k-1}$, equal to $2$ and $4$
	at $k=3$, whence the shift $+6$ and the doubling $2a_2(f)$.
	
	\subsection{Consolidated notation}
	\label{ss:notation}
	
	The catalogue and the trace formula are stated in terms of the following
	weight-$3$ Siegel--paramodular dimensions and their combinations.
	
	\begin{table}[ht]
		\centering
		\renewcommand{\arraystretch}{1.25}
		\begin{tabular}{@{}l p{0.62\linewidth}@{}}
			\toprule
			Symbol & Meaning \\
			\midrule
			$h_2=h_2(p)$ & class number of the principal genus $=\#$ vertices of $G_p$\\
			$h=h(p)$ & supersingular elliptic class number ($=\#$ product factors)\\
			$T_1=T_1(p)$ & type number of the principal genus $=\#$ $\sigma$-orbits\\
			$S_3(\Gamma_0(p)),\,S_3(\Gamma_0'(p)),\,S_3(K(p))$ & weight-$3$ Siegel, Klingen, paramodular cusp spaces\\
			$S_3^{\pm}(K(p))$ & paramodular Atkin--Lehner (Fricke) eigenspaces\\
			$\dim S_k^{\pm}$ & AL-eigenspaces of $\snew_k$, $k\in\{2,4\}$\\
			$\mathrm{ssn}(p)$ & $\dim S_2^+\dim S_4^++\dim S_2^-\dim S_4^-$ (same-sign Yoshida)\\
			$\mathrm{cross}(p)$ & $\dim S_2^-\dim S_4^++\dim S_2^+\dim S_4^-$ (opposite-sign Yoshida)\\
			$\Delta_{\mathrm{III}}(p)$ & $\dim S_3(\Gamma_0(p))-\mathrm{ssn}(p)$ (general-type Siegel meter)\\
			$\delta(p)$ & $\dim S_3(K(p))-\dim S_4^-=\deg N_p$ (number of paramodular non-lifts)\\
			$d(p)$ & $\delta(p)-2\dim S_3^+(K(p))$ (Fricke-signed non-lift defect)\\
			$c_0(\Pi_p)$ & $2a-c-2b$, parahoric-invariant index in Roberts--Schmidt coordinates \cite{RobertsSchmidt}\\
			$\Va_p$ & monic, $\Va_p^2$ the type-$\Va$ block; $\deg\Va_p=\tfrac12\bigl(h_2-1-\dim\snew_4-\mathrm{cross}-\delta\bigr)$\\
			$R(\pi)$ & Fricke involution on the \emph{principal} genus $M_{0,0}(\Upr(p))$; the subscript is omitted throughout\\
			$R_{\mathrm{npg}}(\pi)$ & the same involution on the non-principal genus\\
			\bottomrule
		\end{tabular}
		\caption{Consolidated notation for the weight-$3$ dimension data. The local
			index $c_0$ labels the parahoric-restriction type of $\Pi_p$; its values
			for the relevant types are read from \cite[Table~A.15]{RobertsSchmidt} (as
			in \cite[\S3]{Ibukiyama2018conj}).}
		\label{tab:notation}
	\end{table}
	
	We record the two splittings used repeatedly, both of which follow from the
	definitions in the table:
	\begin{align}
		\dim S_3(\Gamma_0(p)) &= \mathrm{ssn}(p)+\Delta_{\mathrm{III}}(p),\\
		\dim S_3(K(p)) &= \dim S_4^-+\delta(p).
	\end{align}
	The first separates the Siegel cusp space into its same-sign Yoshida part
	and a general-type remainder; the second separates the paramodular cusp
	space into its Saito--Kurokawa part (in bijection with $S_4^-$) and the
	non-lifts counted by $\delta(p)$. Every dimension in
	Table~\ref{tab:notation} is a closed form in $p$ by
	\cite{Ibukiyama2007dim,Ibukiyama-dim}.
	
	Read together, the table encodes one bookkeeping identity: every
	eigensystem is Eisenstein, Saito--Kurokawa, Yoshida, type-$\Va$, or
	general-type, so that
	$h_2=1+\dim\snew_4+\mathrm{cross}(p)+2\deg\Va_p+\delta(p)$, which is the
	definition of $\deg\Va_p$ read backwards. The trace formula refines this
	count by Atkin--Lehner signs: the Yoshida part enters through
	$\dim S_2^-\dim S_4^+-\dim S_2^+\dim S_4^-$, the general-type part through
	the defect $d(p)$, which satisfies $d\le\delta$ with
	$\delta-d=2\dim S_3^+(K(p))$ even; its non-negativity is the sign bias
	proved later and is assumed nowhere. The local index $c_0(\Pi_p)$ enters
	no dimension count; it labels the parahoric type of each eigensystem and
	drives the block multiplicities.
	
	\subsection{Certificate conventions}
	\label{ss:certificate-conventions}
	
	For the finite verification range $11\le p\le149$, the archived data for
	each prime contain an integer matrix $M_p$, automorphism weights $e_i$, an
	involution permutation $\sigma_p$ inducing $R(\pi)$, and the split-side
	block polynomials with their Fricke signs. The verification checks:
	\begin{enumerate}[label=\textup{(Cert\arabic*)},leftmargin=3.7em]
		\item every row of $M_p$ sums to $15$, the Mestre symmetry
		\eqref{eq:mestre} holds, and the weights satisfy the mass identity;
		\item $\sigma_p^2=1$, $\sigma_pM_p=M_p\sigma_p$, and the weights are
		$\sigma_p$-invariant;
		\item the characteristic polynomial of $M_p$ divides exactly by the
		Eisenstein, Saito--Kurokawa, and Yoshida blocks, leaving
		$\Va_p(x)^2N_p(x)$;
		\item for every irreducible factor, the signed isotypic trace
		$\operatorname{Tr}(R\,h_f(M_p))$ matches the predicted signed multiplicity.
	\end{enumerate}
	The fourth check is performed modulo the prime $q=2^{61}-1$. Since the
	signed trace on an isotypic factor is an integer bounded in absolute value
	by the relevant multiplicity, and $2h_2(p)\le2866<q$ throughout the
	certified range, the modular value determines the integer trace once the
	required cofactor inversions are verified factor by factor.
	
	\section{Transfer framework and eigenvalue laws}
	\label{sec:transfer}
	
	This section explains the automorphic framework behind the eigenvalue laws
	of Theorem~\ref{thm:transfer-framework} and the factorization in
	Conjecture~\ref{conj:eigenvalue-sign}. Two questions must be
	distinguished. The transfer and the eigenvalue formulas apply to every
	principal-genus eigensystem that occurs. By contrast, the question of
	which parameters occur in the principal genus, and with what
	multiplicity, is conjectural in general.
	
	\subsection{Block degrees and their non-negativity}
	\label{ss:block-degrees}
	
	The factorization predicted in
	Conjecture~\ref{conj:eigenvalue-sign} has two non-lift degrees. The degree
	$\delta(p)=\deg N_p$ counts paramodular non-lifts; the degree
	$\deg\Va_p$ counts type-$\Va$ pairs. The signed quantity $d(p)$ plays a
	different role: it measures the Fricke imbalance in the general-type
	block. In particular, $d(p)$ is not used in the factorization itself.
	
	\begin{lemma}\label{lem:positivity}
		For every prime $p\ge11$,
		\begin{enumerate}[label=\textup{(\roman*)},leftmargin=2.3em]
			\item $\delta(p)\ge0$;
			\item $\deg\Va_p\ge0$ for $11\le p\le149$, and for every $p\ge11$
			under Conjecture~\ref{conj:eigenvalue-sign}.
		\end{enumerate}
		Statement \textup{(i)} is unconditional; neither statement uses the
		sign of $d(p)$.
	\end{lemma}
	
	\begin{proof}
		For \textup{(i)}, the Saito--Kurokawa part of $S_3(K(p))$ is the space
		of Gritsenko lifts, of dimension $\dim S_4^{-}(\Gamma_0(p))$
		\cite{Gritsenko1995,IPY2013}; it is a subspace of $S_3(K(p))$, so
		$\delta(p)=\dim S_3(K(p))-\dim S_4^{-}\ge0$. In the certified range
		$\delta(p)=0$ for $p\le59$ (Theorem~\ref{thm:certified-spectra}).
		
		For \textup{(ii)}, the certified factorizations of
		Theorem~\ref{thm:certified-spectra} give $\deg\Va_p\ge0$ at every
		prime $11\le p\le149$, and under
		Conjecture~\ref{conj:eigenvalue-sign} the degree $\deg\Va_p$ counts
		the type-$\Va$ parameters occurring in the principal genus, so the
		inequality holds for every~$p$.
	\end{proof}
	
	\begin{remark}
		The bookkeeping definition of $\deg\Va_p$ makes $2\deg\Va_p$ a closed
		form in $p$ (Proposition~\ref{prop:VI}); its non-negativity for
		every prime is exactly the principal-genus multiplicity prediction for
		type~$\Va$, and is established here only in the certified range.
	\end{remark}
	
	\begin{remark}\label{rem:dp-status}
		The inequality $d(p)\ge0$ is a separate statement about the sign distribution
		of paramodular non-lifts. It is not needed in the proofs of the transfer
		theorem, the dimension bookkeeping, or the structural trace identity. It
		holds for every prime: the exact closed formula for $d(p)$ and its
		positivity are proved in the final section (Theorem~\ref{thm:bias}).
	\end{remark}
	
	\subsection{Automorphic decomposition and transfer}
	\label{ss:automorphic-transfer}
	
	Let $G=\GU_2(B)$. The finite-dimensional space
	$M_{0,0}(\Upr(p))$ decomposes into automorphic isotypic components
	\[
	M_{0,0}(\Upr(p))=\bigoplus_{\Pi}V_\Pi,
	\]
	where $\Pi$ runs over automorphic representations of $G$ with trivial
	archimedean component and unramified components at every $q\nmid p$.
	The dimension of $V_\Pi$ is
	\[
	\dim V_\Pi=m(\Pi)\dim\Pi_p^{U_{1,p}},
	\]
	where $m(\Pi)$ is the automorphic multiplicity. The commuting operators
	$B_2(2)$ and $R(\pi)$ preserve each $V_\Pi$.
	
	The Jacquet--Langlands-type transfer from the compact inner form to
	$\GSp_4$ identifies every occurring eigensystem with a split-side Arthur
	parameter. In the present weight and level, only four classes appear:
	Eisenstein, Saito--Kurokawa, Yoshida, and general type; the Borel and
	Klingen CAP classes do not contribute at this weight and level
	\cite[proof of Thm.~9.6]{DPRT}. The transfer and
	multiplicity statements used here are consequences of
	\cite{RoesnerWeissauer2021,vanHoften,DPRT}, organized by Arthur type as in
	\cite{Schmidt2018}.
	
	\begin{proposition}	\label{prop:arthur}
		Let $p\ge11$ and $G=\GU_2(B)$.
		\begin{enumerate}[label=\textup{(\alph*)},leftmargin=2.3em]
			\item Every simultaneous eigensystem of $\{B_2(2),R(\pi)\}$ on
			$M_{0,0}(\Upr(p))$ belongs to an automorphic isotypic component $V_\Pi$
			as above.
			\item Every such $\pi$ transfers to a split-side parameter of exactly one
			of the Arthur types \textup{(Eis)}, \textup{(P)}, \textup{(Y)}, or
			\textup{(G)}. Its Satake parameters agree with those of $\pi$ at every
			prime $q\nmid p$ under the unramified identification
			$G(\Q_q)\simeq\GSp_4(\Q_q)$.
			\item In the spin normalization
			\[
			Q_2(X)=1-\lambda_2X+\cdots+2^6X^4,
			\]
			the eigenvalue of $B_2(2)$ on $V_\Pi$ is the coefficient $\lambda_2$,
			namely
			\[
			\begin{array}{c|c|c}
				\text{type} & \text{spin roots at }2 & \lambda_2\\
				\hline
				\mathrm{Eis} & \{1,2,4,8\} & 15\\
				\mathrm{P} & \{\alpha_g,\beta_g,2,4\} & a_2(g)+6\\
				\mathrm{Y} & \{2\alpha_f,2\beta_f,\alpha_g,\beta_g\}
				& a_2(g)+2a_2(f)\\
				\mathrm{G} & \text{irreducible spin parameter} & a_2(F).
			\end{array}
			\]
		\end{enumerate}
	\end{proposition}
	
	\begin{proof}
		Part \textup{(a)} is the standard automorphic decomposition of the
		finite double-coset space. Part \textup{(b)} follows from the transfer to
		$\GSp_4$, together with the Arthur classification of the split side.
		Part \textup{(c)} follows by taking the sum of the four spin roots. For
		type \textup{(Y)}, the central-character constraint gives the multiplier
		$2$ on the weight-$2$ $\mathrm{GL}_2$-slot. For type \textup{(P)}, the two
		additional roots are $2$ and $4$. Thus the corresponding sums are
		$a_2(g)+2a_2(f)$ and $a_2(g)+6$.
	\end{proof}
	
	\begin{lemma}
		\label{lem:eis}
		The action of $R(\pi)$ on the principal-genus classes is a permutation.
		It fixes the constant functions. Consequently, the Eisenstein eigenvalue
		$15$ has $R(\pi)$-sign $+1$.
	\end{lemma}
	
	\begin{proof}
		The matrix of $R(\pi)$ has exactly one entry equal to $1$ in each row
		\cite[\S3.2]{IbukiyamaKatsura}. Hence it fixes the constant functions.
		Since $B_2(2)$ has row sum $15$, the constant line is the Eisenstein line.
	\end{proof}
	
	\subsection{The local index and the multiplicity table}
	\label{ss:local-index}
	
	For a local representation $\Pi_p$, let $a,c,b$ be the dimensions of its
	fixed-vector spaces under the Klingen, Siegel, and paramodular parahorics,
	respectively. Following \cite[Table~A.15]{RobertsSchmidt} and the
	parahoric-restriction tables of \cite{Roesner2016}, define
	\begin{equation}\label{eq:c0}
		c_0(\Pi_p)=2a-c-2b.
	\end{equation}
	For the relevant Iwahori-spherical types, one has
	$c_0\in\{-1,0,1,2\}$, and $c_0=2$ occurs exactly for type Va.
	
	\begin{table}[ht]
		\centering
		\begin{tabular}{l|c|c}
			Arthur type of $\Pi$ &
			$\dim$ in $M_{0,0}(\Upr(p))$ &
			$\dim$ in $M_{0,0}(U_{\mathrm{npg}}(p))$\\
			\hline
			Eisenstein & $1$ & $1$\\
			Saito--Kurokawa lift of $g\subset\snew_4$ & $1$ &
			$\tfrac12\bigl(1-\eps_p(g)\bigr)$\\
			Yoshida, $\eps_p(f)\eps_p(g)=-1$ & $1$ & $0$\\
			Yoshida, $\eps_p(f)\eps_p(g)=+1$ & $0$ & $0$\\
			general type, local type Va at $p$ \textup{(}$c_0=2$\textup{)} & $2$ & $0$\\
			general type, local $c_0=1$ & $1$ & $1$\\
			types with $c_0\le0$ & $0$ & $0$
		\end{tabular}
		\caption{The two-column multiplicity table at weight $(3,0)$. The
			non-principal column is a theorem. The principal column is the content of
			Ibukiyama's conjectures and is verified by Theorem~\ref{thm:certified-spectra}
			in the range $11\le p\le149$.}
		\label{tab:twocol}
	\end{table}
	
	\begin{table}[ht]
		\centering
		\begin{tabular}{l|c}
			Iwahori-spherical type & $c_0=2a-c-2b$\\
			\hline
			types with a paramodular or Siegel new vector & $\le1$\\
			type Va & $2$
		\end{tabular}
		\caption{The local index from the parahoric fixed-vector dimensions.
			Type Va is the unique relevant type with $c_0=2$.}
		\label{tab:c0}
	\end{table}
	
	The rows are indexed by the global Arthur type of $\Pi$, refined by the
	local type at $p$ in the general-type case. The label type~$\Va$ names the
	local component at $p$ and does not by itself determine the global
	parameter: the Yoshida contribution is counted separately through
	$\mathrm{cross}(p)$, so the type-$\Va$ row consists of general-type
	parameters. The entry $2$ in the principal column is a multiplicity in
	$M_{0,0}(\Upr(p))$, not a dimension of paramodular new vectors: type
	$\Va$ has no paramodular vector at level $p$, its first one occurring at
	level $p^{2}$ \cite[Table~5.2]{RobertsSchmidtTables}, which is the $0$ in
	the non-principal column.
	
	\begin{remark}[The local packet at $p$ and the sign in
		\Cref{conj:eigenvalue-sign}\textup{(iii)}]
		\label{rem:localpacket}
		The $\pm$ splitting asserted for a type-$\Va$ pair is a consequence of
		its multiplicity, not an independent assumption. Write $\xi$ for the
		unramified quadratic character of $\Q_p^\times$. The $L$-parameter of
		type $\Va$ is $\varphi=\varphi_1\oplus\varphi_2$ with
		$\varphi_2=\varphi_1\otimes\xi$, each $\varphi_i$ irreducible of
		dimension $2$; in particular $\varphi_1\not\cong\varphi_2$ and
		$\det\varphi_1=\det\varphi_2$
		\cite[Tables~3.1--3.2]{RobertsSchmidtTables}, and on the split group the
		packet is $\{\Va,\Va^{*}\}$ \cite[\S3]{Schmidt2018}. On the inner form
		$G(\Q_p)=\GU_2(B_p)$ the local Langlands correspondence of
		Gan--Tantono attaches to $\varphi$ a packet parametrized by
		$\widehat{B_\varphi}\smallsetminus\widehat{A_\varphi}$, of size $1$ or
		$2$, and for a parameter of the above shape the size is exactly $2$
		\cite[Main Thm.~(ii), \S7.2]{GanTantono}. Twisting by $\xi\circ\lambda$,
		where $\lambda$ is the similitude character, sends the parameter to
		$\varphi\otimes\xi=\varphi$
		\cite[Main Thm.~(iv)]{GanTantono}, so it permutes this packet. Since
		$\xi$ is unramified and $\lambda$ takes values in $\Z_p^\times$ on
		$\Upr(p)$, the twist acts trivially on $\Upr(p)$-fixed vectors, while
		$\lambda(\pi)=p$ gives $\xi(\lambda(\pi))=-1$: the twist negates
		$R(\pi)$. A member with a one-dimensional $\Upr(p)$-fixed line therefore
		cannot be isomorphic to its own twist, since its $R(\pi)$-eigenvalue
		$r=\pm1$ would satisfy $r=-r$. Hence the twist exchanges the two members
		and their $R(\pi)$-eigenvalues are opposite. The occurrence of both members is proved in
		Section~\ref{sec:occurrence} (Theorem~\ref{thm:occurrence}); what
		remains open in the type-$\Va$ row of Table~\ref{tab:twocol} is the
		local constant of Remark~\ref{rem:local-constant}.
	\end{remark}
	
	The principal column of Table~\ref{tab:twocol} is where the conjectural
	content lies. The existence of the Saito--Kurokawa and Yoshida liftings is
	known, as is the non-principal column. The general-type rows, including the
	type-Va case, are settled in Section~\ref{sec:occurrence} up to the local
	fixed-vector constants of Remarks~\ref{rem:local-constant}
	and~\ref{rem:other-rows}; what remains open in general is the
	principal-column multiplicity statement for the Saito--Kurokawa and Yoshida
	rows.
	
	\subsection{Dimension bookkeeping}
	\label{ss:dimension-bookkeeping}
	
	The block degrees in Section~\ref{sec:statements} are equivalent to the
	known dimension identity of Ibukiyama. We record this equivalence to make
	clear that the degree count is not an additional theorem.
	
	\begin{proposition}
		\label{prop:VI}
		The identity
		\[
		h_2(p)=1+\dim\snew_4+\mathrm{cross}(p)
		+2\deg\Va_p+\delta(p)
		\]
		together with
		\[
		\deg\Va_p
		=\dim S_3(\Gamma_0'(p))-
		\dim S_3(K(p))-\tfrac12\Delta_{\mathrm{III}}(p)-\tfrac12\delta(p)
		\]
		is equivalent, as an identity of closed forms, to
		\begin{equation}\label{eq:ibu-dim}
			2\dim S_3(\Gamma_0'(p))-
			\dim S_3(\Gamma_0(p))-2\dim S_3(K(p))
			=h_2(p)-1-(\dim\snew_2+1)\dim\snew_4,
		\end{equation}
		which is \cite[Thm.~3.1, $j=0$]{Ibukiyama2018conj}.
	\end{proposition}
	
	\begin{proof}
		Substitute
		\[
		\dim S_3(\Gamma_0(p))=\mathrm{ssn}(p)+\Delta_{\mathrm{III}}(p),
		\qquad
		\dim S_3(K(p))=\dim S_4^-+\delta(p),
		\]
		and
		\[
		\dim\snew_2\dim\snew_4
		=\mathrm{ssn}(p)+\mathrm{cross}(p)
		\]
		into \eqref{eq:ibu-dim}. The same-sign term cancels, and rearrangement
		gives exactly the asserted block-degree formula.
	\end{proof}
	
	\begin{remark}
		\label{rem:transfer-scope}
		The transfer theorem and Proposition~\ref{prop:arthur} determine the
		$B_2(2)$-eigenvalue attached to an occurring parameter. They do not
		determine which split-side parameters occur in the principal genus, nor
		their individual multiplicities. For the general-type rows those assertions are settled in
		Section~\ref{sec:occurrence}; for the Saito--Kurokawa and Yoshida rows
		they belong to Conjecture~\ref{conj:eigenvalue-sign} and are
		established in this paper only in the finite range of
		Theorem~\ref{thm:certified-spectra}.
	\end{remark}
	
	\section{The occurrence theorem for the type-Va block}
	\label{sec:occurrence}
	
	This section proves the occurrence half of
	Conjecture~\ref{conj:eigenvalue-sign}\textup{(iii)}. The input is the
	global lifting theorem of R\"osner and Weissauer for inner forms of
	$\GSp_4$ that are anisotropic at the archimedean place
	\cite[Thm.~11.4]{RoesnerWeissauer2021}. Applied to $G=\GU_2(B)$ at
	trivial weight, it shows that the weak packet of a general-type
	automorphic representation is the full product of its local $L$-packets
	and that every member occurs with multiplicity one. The $\pm$ splitting
	of Remark~\ref{rem:localpacket} then holds unconditionally, so the
	type-Va contribution to $\charp B_2(2)$ is an exact square split evenly
	by $R(\pi)$ at every prime. The open content of the type-Va row of
	Table~\ref{tab:twocol} is thereby reduced to a single local constant
	(Remark~\ref{rem:local-constant}).
	
	\subsection{Applicability of the lifting theorem}
	\label{ss:rw-applicability}
	
	\begin{lemma}
		\label{lem:rw-applicable}
		The group $G=\GU_2(B)$ satisfies the standing hypotheses of
		\cite[\S1]{RoesnerWeissauer2021}: its center is
		$Z\cong\mathbb{G}_m$, so $H^1(\Q,Z)=1$, and its derived group is an
		inner form of $\Sp_4$, hence simply connected. The real group
		$G(\R)$ is anisotropic modulo center. Consequently Conjecture~7.5 of
		\cite{RoesnerWeissauer2021} holds for $G$, the virtual
		multiplicities of \cite[\S9]{RoesnerWeissauer2021} are cuspidal
		multiplicities, and \cite[Thm.~11.4]{RoesnerWeissauer2021} applies
		to every automorphic representation of $G(\mathbb{A})$ of general
		type with trivial archimedean component. No regularity hypothesis on
		the weight is required.
	\end{lemma}
	
	\begin{proof}
		The center consists of the scalars $z\cdot1_2$ with
		$z\in\mathbb{G}_m$, and $H^1(\Q,\mathbb{G}_m)=1$ by Hilbert's
		Theorem~90. The derived group is the reduced-norm-one subgroup, an
		inner form of the simply connected group $\Sp_4$. Since $B$ ramifies
		at $\infty$, the Hermitian form is definite and $G(\R)$ is compact
		modulo center; the associated double coset spaces are finite, so
		their cohomology is concentrated in degree $0$ and Conjecture~7.5 of
		\cite{RoesnerWeissauer2021} holds, with $d(G)=1$ and virtual equal
		to cuspidal multiplicities; this is the algebraic modular form case
		treated after \cite[Cor.~7.4]{RoesnerWeissauer2021}, and the
		coefficient normalization is \cite[Lem.~9.1,
		Prop.~9.2]{RoesnerWeissauer2021}. Theorem~11.4 of
		\cite{RoesnerWeissauer2021} is stated for inner forms anisotropic at
		the archimedean place with no regularity condition on the
		coefficient system; the regular-weight statement for the remaining
		inner forms is their Theorem~11.5 and is not used here. The same
		specialization, at scalar weight $3$ on the definite form, is
		applied in the proof of \cite[Prop.~12.3]{RoesnerWeissauer2021},
		where an additional argument is needed only when the inner form is
		split at the archimedean place.
	\end{proof}
	
	\subsection{Occurrence and sign splitting}
	\label{ss:occurrence-theorem}
	
	Write $U_{1,p}\subset G(\Q_p)$ for the principal-genus parahoric, so
	that $\Upr(p)$ has component $U_{1,p}$ at $p$ and hyperspecial
	components elsewhere, and let $\xi$ and $\lambda$ be as in
	Remark~\ref{rem:localpacket}.
	

	\begin{theorem}[Occurrence and sign splitting for type Va]
		\label{thm:occurrence}
		Let $p\ge11$ and let $\Pi$ be an irreducible automorphic
		representation of $G(\mathbb{A})$ of general type contributing to
		$M_{0,0}(\Upr(p))$ whose local component $\Pi_p$ lies in the
		$L$-packet attached by Gan--Tantono to a parameter of type $\Va$;
		the packet is $\{\Pi_p^+,\Pi_p^-\}$ with
		$\Pi_p^-\cong(\xi\circ\lambda)\otimes\Pi_p^+$, denoted
		$\{\Va_G,\chi_0\otimes\Va_G\}$ in
		\cite[Table~3]{RoesnerWeissauer2021}. Write $\Pi^{\pm}$ for the
		representations with the same components as $\Pi$ away from $p$ and
		with $p$-component $\Pi_p^{\pm}$, and
		$V[\Pi]\subset M_{0,0}(\Upr(p))$ for the sum of the corresponding
		isotypic subspaces. Then:
		\begin{enumerate}[label=\textup{(\alph*)},leftmargin=2.3em]
			\item Both $\Pi^+$ and $\Pi^-$ are automorphic, each with
			multiplicity one, and they exhaust the weak equivalence class of
			$\Pi$.
			\item $V[\Pi]=(\Pi_p^+)^{U_{1,p}}\oplus(\Pi_p^-)^{U_{1,p}}$. The
			two summands are $R(\pi)$-stable of equal dimension
			$d_p(\Pi):=\dim(\Pi_p^+)^{U_{1,p}}\ge1$, and the multiplicities
			of the $R(\pi)$-eigenvalues $+1$ and $-1$ on the two summands
			are interchanged. In particular $\dim V[\Pi]=2\,d_p(\Pi)$, the
			eigenvalues $+1$ and $-1$ of $R(\pi)$ on $V[\Pi]$ each occur
			with multiplicity $d_p(\Pi)$, and
			$\Tr\bigl(R(\pi)\mid V[\Pi]\bigr)=0$.
			\item Every $T(n)$ with $p\nmid n$ acts on $V[\Pi]$ by a single
			scalar; in particular the eigensystems in $V[\Pi]$ agree at
			every such $T(n)$ and are separated only by $R(\pi)$. The
			contribution of the weak packet of $\Pi$ to $\charp B_2(2)$ is
			$(x-\lambda_2(\Pi))^{2d_p(\Pi)}$ for the common
			$B_2(2)$-eigenvalue $\lambda_2(\Pi)$ of
			Proposition~\ref{prop:arthur}, and the total type-Va
			contribution to $\charp B_2(2)$ is an exact square on which
			$R(\pi)$ has trace zero and equal $\pm1$ eigenvalue
			multiplicities.
			\item Conversely, every cuspidal automorphic representation
			$\Pi'$ of $\GSp_4(\mathbb{A})$ of general type, cohomological of
			weight $(3,0)$, unramified outside $p$, whose $p$-component lies
			in the packet $\{\Va,\Va^{*}\}$, is a lifting from $G$ of the
			above kind, and the induced correspondence of parameters between
			such $\Pi'$ and the weak packets in \textup{(a)} is a bijection.
		\end{enumerate}
	\end{theorem}

	\begin{proof}
		\textup{(a)} By Lemma~\ref{lem:rw-applicable},
		\cite[Thm.~11.4]{RoesnerWeissauer2021} applies to $\Pi$.
		Assertion~1 of that theorem gives multiplicity one for every
		automorphic representation weakly equivalent to $\Pi$, and
		assertion~2 identifies the set of these representations with the
		product of the local $L$-packets of the components of $\Pi$, the
		local packets corresponding to the local Langlands parameters. The
		archimedean packet is the singleton consisting of the trivial
		representation, the packet at every $q\ne p$ is the singleton
		containing the unramified member, and the packet at $p$ is
		$\{\Pi_p^+,\Pi_p^-\}$, of cardinality two
		\cite[Main Thm.~(ii), \S7.2]{GanTantono}. Hence the weak
		equivalence class of $\Pi$ is $\{\Pi^+,\Pi^-\}$, both automorphic
		with multiplicity one. Corollary~11.6 of
		\cite{RoesnerWeissauer2021}, which lists both $\Va_G$ and
		$\chi_0\otimes\Va_G$ among the local types of general-type
		representations of inner forms, records the same conclusion.
		
		\textup{(b)} Multiplicity one gives
		$V[\Pi]=(\Pi_p^+)^{U_{1,p}}\oplus(\Pi_p^-)^{U_{1,p}}$, with
		$B_2(2)$ and $R(\pi)$ acting through $\Pi_p^{\pm}$ on the summands;
		$R(\pi)$ acts through a local element $\omega_p$ normalizing
		$U_{1,p}$ with $\lambda(\omega_p)=p$, as in
		Remark~\ref{rem:localpacket}. Fix an isomorphism
		$\iota\colon(\xi\circ\lambda)\otimes\Pi_p^+\to\Pi_p^-$. Since $\xi$
		is unramified and $\lambda(U_{1,p})\subset\Z_p^\times$, the twist
		is trivial on $U_{1,p}$, so $\iota$ restricts to an isomorphism
		$(\Pi_p^+)^{U_{1,p}}\to(\Pi_p^-)^{U_{1,p}}$; in particular the two
		dimensions agree, and they are nonzero because $\Pi$ contributes to
		$M_{0,0}(\Upr(p))$. Under $\iota$ the action of $\omega_p$ on
		$(\Pi_p^-)^{U_{1,p}}$ corresponds to
		$\xi(\lambda(\omega_p))=\xi(p)=-1$ times its action on
		$(\Pi_p^+)^{U_{1,p}}$. Since $R(\pi)^2=\mathrm{id}$, the
		eigenvalues are $\pm1$ with multiplicities $(a,b)$ on the first
		summand and $(b,a)$ on the second, whence the balanced spectrum and
		the vanishing trace.
		
		\textup{(c)} The members of a weak packet share Satake parameters
		at every $q\nmid p$, so every $T(n)$ with $p\nmid n$, and in
		particular $B_2(2)$, acts on $V[\Pi]$ by the scalar determined by
		Proposition~\ref{prop:arthur}. The characteristic polynomial of
		$B_2(2)$ on $V[\Pi]$ is therefore
		$(x-\lambda_2(\Pi))^{2d_p(\Pi)}$, and the product over the weak
		packets in question is a square; $R(\pi)$ is balanced on each
		factor by \textup{(b)}.
		
		\textup{(d)} By \cite[Thm.~11.4(4)]{RoesnerWeissauer2021}, with the
		archimedean weight matching stated in
		\cite[Cor.~11.7]{RoesnerWeissauer2021}, a general-type cohomological
		$\Pi'$ is a lifting from $G$ if and only if at every finite place
		where $G$ is not split the local $L$-packet of $\Pi'$ is not fully
		Klingen induced. Here the only such place is $p$, and the packet
		$\{\Va,\Va^{*}\}$ contains the supercuspidal $\Va^{*}$
		\cite[\S3]{Schmidt2018}, so it is not the semisimplification of any
		parabolically induced representation; equivalently, type $\Va$ is
		not among the excluded types of
		\cite[Cor.~11.7]{RoesnerWeissauer2021}. The bijectivity of the
		correspondence of parameters follows from assertion~2 of
		\cite[Thm.~11.4]{RoesnerWeissauer2021} together with the fact that
		the spherical data determine the split-side weak packet
		\cite[Prop.~10.1(6)]{RoesnerWeissauer2021}.
	\end{proof}

	
	\begin{corollary}
		\label{cor:conj-iii}
		For every prime $p\ge11$, the contribution to $\charp B_2(2)$ of
		the general-type eigensystems whose local component at $p$ is of
		type $\Va$ is an exact square, and $R(\pi)$ splits it into $+1$ and
		$-1$ eigenspaces of equal dimension; its $R(\pi)$-trace vanishes.
		This establishes the square and sign assertions of
		Conjecture~\ref{conj:eigenvalue-sign}\textup{(iii)}. Identifying
		this square with the residual factor $\Va_p(x)^2$ of
		\eqref{eq:C1-factorization} amounts to the remaining rows of
		Table~\ref{tab:twocol}, that is, to the Saito--Kurokawa and Yoshida
		entries of the principal column together with the local
		fixed-vector statements of Remarks~\ref{rem:local-constant}
		and~\ref{rem:other-rows}.
	\end{corollary}
	
	\begin{remark}\label{rem:local-constant}
		Theorem~\ref{thm:occurrence} reduces the type-Va row of
		Table~\ref{tab:twocol} to the single local quantity
		$d_p(\Pi)=\dim(\Pi_p^+)^{U_{1,p}}$, the dimension of the
		$U_{1,p}$-fixed vectors of a packet member; the entry $2$ of the
		row is equivalent to $d_p(\Pi)=1$ for every occurring parameter. At
		$p=19$ the certified factorization
		(Theorem~\ref{thm:certified-spectra}) exhibits a type-Va block of
		dimension two; since each occurring parameter contributes an even
		dimension at least two, there is exactly one type-Va parameter at
		$p=19$ and $d_{19}=1$. For general $p$ the equality $d_p=1$ is a
		finite local computation in the parahoric restriction of the
		Iwahori-spherical representations of $\GU_2(B_p)$; the split-side
		analogue is carried out in \cite{Roesner2016}, and we are not aware
		of a published counterpart for the inner form.
	\end{remark}
	
	\begin{remark}\label{rem:other-rows}
		The argument of Theorem~\ref{thm:occurrence}\textup{(a)} applies to
		every general-type row of Table~\ref{tab:twocol}: by
		\cite[Thm.~11.4]{RoesnerWeissauer2021}, each occurring general-type
		parameter contributes, with multiplicity one, every member of its
		weak packet, so each principal-column entry in a general-type row
		equals the sum of the $U_{1,p}$-fixed dimensions over the local
		packet at $p$. The entries $1$ for $c_0=1$ (local type
		$\mathrm{IIa}_G$; the packet is a singleton, cf.\
		\cite[\S3]{Schmidt2018} and \cite[Table~3]{RoesnerWeissauer2021})
		and $0$ for $c_0\le0$ are therefore equivalent to local statements
		of the same kind as Remark~\ref{rem:local-constant}.
		Corollary~11.6 of \cite{RoesnerWeissauer2021} bounds the possible
		local types at $p$ of occurring general-type representations.
	\end{remark}
	
	R\"osner and Weissauer prove the split-side input to
	\cite[Thm.~11.4]{RoesnerWeissauer2021}, the packet structure of the
	general-type cohomological spectrum of $\GSp_4$, independently of the
	announced classification results
	(\cite[\S10]{RoesnerWeissauer2021}); the proof of the lifting theorem
	proceeds through the cohomological trace formula and the local
	character identities of Chan and Gan \cite{ChanGan2015}. The route
	through Arthur's formalism, and with it the transfer-factor comparison
	named as its prerequisite by Gee and Ta\"ibi
	\cite[Rem.~7.4.8]{GeeTaibi}, is therefore not needed for the present
	application; see \S\ref{sec:concl} for what remains open.
	
	
	\section{The structural trace identity}
	\label{sec:trace}
	
	This section proves Theorem~\ref{thm:structural-trace}. The trace of the
	Atkin--Lehner involution on each genus has a closed evaluation due to
	Ibukiyama. Comparing the principal and non-principal formulas term by term
	produces a simple genus-difference identity. The non-principal-genus
	correspondence then turns this identity into the structural formula for the
	principal genus.
	
	\subsection{Conventions for the trace formulas}
	\label{ss:trace-conventions}
	
	Throughout, $h(\sqrt{-m})$ denotes the class number of the imaginary
	quadratic field $\Q(\sqrt{-m})$. Thus its discriminant is the fundamental
	discriminant attached to $-m$; this convention differs from an
	order-discriminant convention in which $h(-D)$ is set to zero unless
	$-D\equiv0,1\pmod4$.
	
	Let
	\[
	s:=\kro{2}{p},\qquad e_3:=\kro{p}{3},
	\]
	and define
	\begin{equation}\label{eq:nu-principal}
		\nu(p):=
		\begin{cases}
			h(\sqrt{-p}),&p\equiv1\pmod4,\\
			(3-s)h(\sqrt{-p}),&p\equiv3\pmod4.
		\end{cases}
	\end{equation}
	Equivalently, in discriminant notation,
	\[
	\nu(p)=
	\begin{cases}
		h(-4p),&p\equiv1\pmod4,\\
		h(-4p)+h(-p),&p\equiv3\pmod4.
	\end{cases}
	\]
	For example, $\nu(11)=h(-44)+h(-11)=4$ and
	$\nu(23)=h(-92)+h(-23)=6$.
	
	The notation $B_{2,\chi}$ denotes the second generalized Bernoulli number
	for the real quadratic character attached to $\Q(\sqrt p)$, with the
	normalization of \cite{IbukiyamaKatsura,Ibukiyama2019quinary}.
	
	\subsection{Comparison of the two genus traces}
	\label{ss:trace-comparison}
	
	\begin{proposition}
		\label{prop:genus-trace-comparison}
		For every prime $p\ge7$,
		\begin{equation}\label{eq:genus-trace-difference}
			\Tr\!\left(R(\pi)\mid M_{0,0}(\Upr(p))\right)
			-
			\Tr\!\left(R_{\mathrm{npg}}(\pi)\mid M_{0,0}(U_{\mathrm{npg}}(p))\right)
			=
			\nu(p)\frac{p-e_3}{12}.
		\end{equation}
	\end{proposition}
	
	\begin{proof}
		For $p\equiv3\pmod4$, \cite[Rem.~3]{IbukiyamaKatsura} gives
		\begin{align}
			\Tr R(\pi)
			={}&\frac{1}{2^5\cdot3}B_{2,\chi}
			+\frac18h(\sqrt{-2p})
			+\frac1{12}h(\sqrt{-3p}) \notag\\
			&+\left\{
			\frac{(p-1)(9-4s)}{48}
			+\frac{p-s}{16}
			+\frac{(1-e_3)(3-s)}{12}
			\right\}h(\sqrt{-p}).
			\label{eq:principal-trace-3mod4}
		\end{align}
		For $p\equiv1\pmod4$, the same source gives
		\begin{align}
			\Tr R(\pi)
			={}&\frac{9-2s}{2^5\cdot3}B_{2,\chi}
			+\frac{4p-1}{48}h(\sqrt{-p})
			+\frac18h(\sqrt{-2p}) \notag\\
			&+\frac{3+\kro{-2}{p}}{12}h(\sqrt{-3p})
			+\frac{1-e_3}{12}h(\sqrt{-p}).
			\label{eq:principal-trace-1mod4}
		\end{align}
		
		For the non-principal genus, \cite[Thm.~5.2]{Ibukiyama2019quinary} gives,
		when $p\equiv1\pmod4$,
		\begin{equation}
			\Tr R_{\mathrm{npg}}(\pi)
			=\frac{9-2s}{2^5\cdot3}B_{2,\chi}
			+\frac1{16}h(\sqrt{-p})
			+\frac18h(\sqrt{-2p})
			+\frac{3+s}{12}h(\sqrt{-3p}),
			\label{eq:TOH1}
		\end{equation}
		and, when $p\equiv3\pmod4$,
		\begin{equation}
			\Tr R_{\mathrm{npg}}(\pi)
			=\frac{1}{2^5\cdot3}B_{2,\chi}
			+\frac{1-s}{16}h(\sqrt{-p})
			+\frac18h(\sqrt{-2p})
			+\frac1{12}h(\sqrt{-3p}).
			\label{eq:TOH3}
		\end{equation}
		
		The coefficient comparison is summarized in
		Table~\ref{tab:trace-coefficients}. The $B_{2,\chi}$ and
		$h(\sqrt{-2p})$ coefficients agree in both residue classes. For
		$p\equiv1\pmod4$,
		\[
		\kro{-2}{p}=\kro{-1}{p}\kro{2}{p}=s,
		\]
		so the $h(\sqrt{-3p})$ coefficients agree; for $p\equiv3\pmod4$ they
		are both $1/12$. Thus only the $h(\sqrt{-p})$ coefficient remains.
		
		If $p\equiv1\pmod4$, its difference is
		\[
		\frac{4p-1}{48}+\frac{1-e_3}{12}-\frac1{16}
		=\frac{p-e_3}{12}.
		\]
		If $p\equiv3\pmod4$, it is
		\begin{align*}
			&\frac{(p-1)(9-4s)}{48}
			+\frac{p-s}{16}
			+\frac{(1-e_3)(3-s)}{12}
			-\frac{1-s}{16}\\
			&\hspace{4em}=
			\frac{(3-s)(p-e_3)}{12}.
		\end{align*}
		The definition \eqref{eq:nu-principal} gives
		\eqref{eq:genus-trace-difference} in both cases.
	\end{proof}
	
	\begin{table}[ht]
		\centering
		\begin{tabular}{c|c|c}
			Term & $p\equiv1\pmod4$ & $p\equiv3\pmod4$ \\
			\hline
			$B_{2,\chi}$ & equal & equal \\
			$h(\sqrt{-2p})$ & equal & equal \\
			$h(\sqrt{-3p})$ & equal, since $\left(\frac{-2}{p}\right)=\left(\frac2p\right)$ & equal \\
			$h(\sqrt{-p})$ & $(p-e_3)/12$ & $(3-s)(p-e_3)/12$ \\
		\end{tabular}
		\caption{Difference ``principal trace minus non-principal trace'' in the
			closed trace formulas. Only the $h(\sqrt{-p})$ coefficient survives.}
		\label{tab:trace-coefficients}
	\end{table}
	
	\begin{remark}
		\label{rem:global-orbital}
		Both trace formulas arise from the Eichler--Selberg trace formula for
		$\GU_2(B)$, with test functions that differ only at $p$, where the two
		genera correspond to non-conjugate maximal parahorics. The coefficient
		comparison above may be interpreted as equality of the aggregated
		contributions represented by the order-$3$ and order-$4$ torsion terms.
		We do not claim here a direct termwise local equality of twisted orbital
		integrals. Such a statement requires fixed Haar measures, explicit local
		representatives, and a local optimal-embedding calculation.
	\end{remark}
	
	\subsection{Elliptic Atkin--Lehner input}
	\label{ss:elliptic-trace-input}
	
	Write
	\[
	d_k=\dim\snew_k,\qquad
	a_k=\dim S_k^+,
	\qquad b_k=\dim S_k^-
	\qquad(k=2,4).
	\]
	The following identities are used to convert the genus comparison into the
	principal-genus trace formula.
	
	\begin{lemma}	\label{lem:elliptic-input}
		For every prime $p\ge7$,
		\begin{enumerate}[label=\textup{(\alph*)},leftmargin=2.3em]
			\item $a_2-b_2=1-\nu(p)/2$;
			\item $a_4-b_4=\nu(p)/2$;
			\item $d_2+d_4+1=(p-e_3)/3$.
		\end{enumerate}
	\end{lemma}
	
	\begin{proof}
		Part \textup{(a)} follows from the fixed points of the Fricke involution
		on $X_0(p)$ and Riemann--Hurwitz \cite[\S2]{Ogg1974}. Part \textup{(b)} is Yamauchi's trace
		formula in weight $4$ \cite{Yamauchi1973}. Part \textup{(c)} follows by combining the standard
		dimension formulas for weights $2$ and $4$ with quadratic reciprocity; see
		\cite{AtkinLehner,Ibukiyama-dim}.
	\end{proof}
	
	\subsection{Proof of the structural trace identity}
	\label{ss:proof-structural-trace}
	
	The non-principal-genus correspondence \cite{DPRT,Ibukiyama-dim} gives
	\begin{equation}\label{eq:npg-dimension-trace}
		H_{\mathrm{npg}}(p)=1+\dim S_3(K(p)),
		\qquad
		\Tr R_{\mathrm{npg}}(\pi)
		=H_{\mathrm{npg}}(p)-2\dim S_3^+(K(p)).
	\end{equation}
	Here $S_3^\pm(K(p))$ are the eigenspaces of the paramodular Fricke
	involution; the correspondence matches this involution with
	$R_{\mathrm{npg}}(\pi)$ itself, and the signed identity
	$\dim S_3^{+}(K(p))=\tfrac12\bigl(H_{\mathrm{npg}}(p)-\Tr R_{\mathrm{npg}}(\pi)\bigr)$
	is the one used in \cite[proof of Prop.~8.2]{Ibukiyama-dim}, so the
	second identity carries no free sign normalization.
	Since
	\[
	\delta(p)=\dim S_3(K(p))-\dim S_4^-,
	\qquad
	d(p)=\delta(p)-2\dim S_3^+(K(p)),
	\]
	we obtain
	\begin{equation}\label{eq:npg-trace-d}
		\Tr R_{\mathrm{npg}}(\pi)=1+\dim S_4^-+d(p).
	\end{equation}
	
	\begin{proof}[Proof of Theorem~\ref{thm:structural-trace}]
		By Proposition~\ref{prop:genus-trace-comparison} and
		\eqref{eq:npg-trace-d}, it is enough to prove
		\begin{equation}\label{eq:lift-comparison-target}
			L_0(p)-1-\dim S_4^-=
			\nu(p)\frac{p-e_3}{12}.
		\end{equation}
		Using $d_k=a_k+b_k$, the left side is
		\[
		L_0(p)-1-b_4=b_2a_4-a_2b_4.
		\]
		By Lemma~\ref{lem:elliptic-input},
		\[
		b_2a_4-a_2b_4
		=\frac14(d_2-a_2+b_2)(d_4+a_4-b_4)
		-\frac14(d_2+a_2-b_2)(d_4-a_4+b_4).
		\]
		Substituting parts \textup{(a)}--\textup{(c)} of that lemma and simplifying
		gives exactly the right-hand side of \eqref{eq:lift-comparison-target}.
		Therefore
		\[
		\Tr R(\pi)
		=\Tr R_{\mathrm{npg}}(\pi)+\nu(p)\frac{p-e_3}{12}
		=L_0(p)+d(p),
		\]
		as required.
	\end{proof}
	
	\begin{remark}	\label{rem:trace-consistency}
		If Conjecture~\ref{conj:eigenvalue-sign} holds, the Eisenstein and
		Saito--Kurokawa blocks contribute $1+\dim\snew_4$ to the trace. The
		opposite-sign Yoshida blocks contribute
		\[
		\dim S_2^-\dim S_4^+-\dim S_2^+\dim S_4^-.
		\]
		Every type-$\Va$ pair contributes zero, and the general-type contribution
		is $d(p)$. Thus Theorem~\ref{thm:structural-trace} proves the signed total
		predicted by C1 for every prime. It does not, by itself, determine the
		individual multiplicities or eigenvalue assignments of C1.
	\end{remark}
	
	\subsection{Consequences for type numbers}
	\label{ss:type-number-trace}
	
	The involution $R(\pi)$ permutes the principal-genus classes. Hence
	\[
	\Tr R(\pi)=2T_1(p)-h_2(p),
	\]
	where $T_1(p)$ is the principal-genus type number
	\cite{IbukiyamaKatsura,Ibukiyama2018type}.
	
	\begin{corollary}\label{cor:typenum}
		For every prime $p\ge7$,
		\begin{equation}\label{eq:type-number-paramodular}
			T_1(p)=\frac12\left(h_2(p)+L_0(p)+d(p)\right).
		\end{equation}
		Equivalently,
		\[
		T_1(p)=\frac12\left(
		h_2(p)+1+\dim\snew_4
		+\dim S_2^-\dim S_4^+
		-\dim S_2^+\dim S_4^-
		+d(p)
		\right).
		\]
	\end{corollary}
	
	\begin{proof}
		Combine $\Tr R(\pi)=2T_1(p)-h_2(p)$ with
		Theorem~\ref{thm:structural-trace}.
	\end{proof}
	
	The value of $T_1(p)$ itself is classical. The contribution of
	Corollary~\ref{cor:typenum} is its expression in paramodular dimension data.
	
	\begin{example}\label{ex:trace-p11}
		At $p=11$, one has $h_2=5$, $\dim\snew_4=2$,
		$\dim S_2^-=1$, $\dim S_4^+=2$, and $d(11)=0$. Thus
		\[
		\Tr R(\pi)=1+2+1\cdot2=5.
		\]
		This agrees with the fact that all five principal-genus classes are fixed
		by the involution, and gives $T_1(11)=5$.
	\end{example}
	
	\section{Certified computation and the proof of Theorem~\ref{thm:certified-spectra}}
	\label{sec:certificates}
	
	This section proves Theorem~\ref{thm:certified-spectra}. For each prime $11\le p\le149$, the archived artifact supplies a finite list of integer data, and a small number of exact identities among these data imply the full conclusion of Conjecture~\ref{conj:eigenvalue-sign} at $p$. The format follows the
	tradition of succinct certificates and certificate-checked computer proof
	\cite{Pratt1975,AtkinMorain1993,RSST1997}. The section is organized in
	three layers: the construction of the geometric data, the deductive step
	that converts certificates into the spectral conclusion, and the
	reproducibility protocol that makes the computation independently
	checkable.
	
	\subsection{Overview of the logical scheme}
	\label{ss:cert-overview}
	
	The proof separates into a deductive proposition (Proposition~\ref{prop:reduction}) and a finite computation, replayed from the archived data in exact arithmetic.
	
	\begin{enumerate}[label=\textup{(\arabic*)},leftmargin=2.4em]
		\item \textbf{Geometric data.} For each prime, the Richelot engine produces
		an integer out-edge matrix $M_p$, a permutation $\sigma_p$ inducing
		$R(\pi)$, and vertex weights $e_i$, from the geometry alone
		(\S\ref{ss:engine}, \S\ref{ss:extension}).
		\item \textbf{Split-side data.} Two separate pipelines supply the weight-$2$
		and weight-$4$ newform orbit polynomials with their Atkin--Lehner signs
		(\S\ref{ss:split-pipelines}).
		\item \textbf{Certificates.} Finitely many exact identities among these
		data are verified in exact arithmetic (\S\ref{ss:finitecert}).
		\item \textbf{Conclusion.} A deductive proposition
		(Proposition~\ref{prop:reduction}) converts the verified identities into
		the simultaneous eigenspace and sign decomposition asserted by
		Conjecture~\ref{conj:eigenvalue-sign} at $p$.
	\end{enumerate}
	
	External comparisons with paramodular databases are reported in
	\S\ref{ss:external} as corroboration; they are not inputs to the proof.
	
	\subsection{The geometric engine}
	\label{ss:engine}
	
	The engine constructs the out-edge matrix $M_p$ from Igusa invariants
	alone. Fix a superspecial Jacobian vertex $A\simeq\Jac C$ with
	$C:y^2=f(x)$, $\deg f=6$, over a field of definition $\F_{p^2}$. Each of
	the fifteen Richelot kernels at $A$ corresponds to a partition of the six
	roots of $f$ into three unordered pairs, equivalently to a factorization
	$f=c\,q_1q_2q_3$ into quadratics over the splitting field
	\cite[\S8]{JordanZaytman}. For each partition the codomain is determined by
	the Richelot determinant
	\[
	\Delta_R=\det\begin{pmatrix}a_0&a_1&a_2\\ b_0&b_1&b_2\\ c_0&c_1&c_2
	\end{pmatrix},
	\qquad q_i=a_ix^2+b_ix+c_i:
	\]
	if $\Delta_R\ne0$ the target is the Jacobian of
	$\Delta_R y^2=[q_2,q_3][q_3,q_1][q_1,q_2]$ with
	$[q_i,q_j]=q_i'q_j-q_iq_j'$, and if $\Delta_R=0$ the target is a product
	of elliptic curves. Two targets coincide as principally polarized abelian
	surfaces precisely when their Igusa--Clebsch invariants agree in weighted
	projective space \cite{Igusa1960}; this is the test used to locate the
	target vertex.
	
	The engine evaluates the bracket polynomials $[q_i,q_j]$ and the
	invariants through subresultants of $f$ and its partition factors, without
	forming the derivative $f'$ and without passing to a splitting field \cite{Dang2025}. The
	subresultant chain computes the required resultants and brackets over the
	ground ring directly, so the codomain invariants are obtained without the
	degree drops that a derivative-based implementation incurs at ramified
	points. All arithmetic is exact. Product vertices $E_a\times E_b$ are
	handled by the companion incidence of \cite{KatsuraTakashima}, their
	fifteen kernels distributed among Jacobian and product targets according to
	the isotropic-subspace combinatorics. The result is the integer matrix
	$M=(M_{ij})$ with $M_{ij}$ the number of kernels at $i$ with target $j$,
	and $\sum_jM_{ij}=15$ for every $i$.
	
	\subsection{Self-weighting}
	\label{ss:selfweight}
	
	The operator $B_2(2)$ is the weighted adjacency operator
	$S=D^{-1/2}MD^{1/2}$, $D=\diag(e_i)$, not the raw matrix $M$. The weights
	$e_i=\#\Aut(A_i,\lambda_i)$ are recovered from $M$ itself rather than
	supplied to it. By Lemma~\ref{lem:vertices}(b), the product-vertex weights
	are $8w_a^2$ and $4w_aw_b$ in terms of the elliptic automorphism numbers
	$w_a\in\{1,2,3\}$, and the diagonal product weights $e_{\{a,a\}}=8w_a^2$
	determine the $w_a$; the Eichler mass $\sum_a1/w_a=(p-1)/12$ of
	\eqref{eq:eichler-mass} fixes their global scale. The Jacobian-vertex
	weights are then pinned by the Mestre symmetry \eqref{eq:mestre}: along
	any two-way edge $M_{ij},M_{ji}>0$ the relation $e_jM_{ij}=e_iM_{ji}$
	propagates a known weight to its neighbour, and the Jacobian balance graph
	is connected at every prime in the certified range, so a single product
	anchor determines all remaining weights. The propagation is
	over-determined: every two-way ordered edge imposes one independent
	instance of \eqref{eq:mestre}, and all hold simultaneously as identities
	in $\Q$. A single transcription error in any $M_{ij}$ would violate the
	symmetry at some edge and be detected.
	
	\subsection{Three structural certificates}
	\label{ss:pillars}
	
	Each output matrix is required to pass three internal checks, all in exact
	arithmetic and independent of the automorphic side.
	
	\begin{enumerate}[label=\textup{(P\arabic*)},leftmargin=2.6em]
		\item \emph{Combinatorial validity.} Every row of $M$ sums to $15$; the
		Mestre symmetry \eqref{eq:mestre} holds for the self-weighting weights;
		and the recovered weights satisfy the Eichler mass \eqref{eq:eichler-mass}.
		This certifies that $M$ is the out-edge matrix of a $15$-regular Richelot
		graph with the correct total mass.
		\item \emph{Spectral integrality.} The characteristic polynomial of $M$,
		recomputed from the integer entries, factors over $\Z$ with the Perron
		factor $(x-15)$ simple (Lemma~\ref{lem:vertices}(d)). This is checked
		independently of the system that produced $M$, so a transcription error in
		any entry is caught by a mismatch in $\charp(M)$.
		\item \emph{Local cross-check.} The Jacobian/product incidence and the
		edge multiplicities agree with the independent enumerations of
		\cite{KatsuraTakashima,FloritSmith}, which count superspecial Richelot
		isogenies by a different method.
	\end{enumerate}
	
	\subsection{Extension to every prime $p\le149$}
	\label{ss:extension}
	
	The engine of \S\ref{ss:engine} assumes the six Weierstrass points of each
	Jacobian vertex rational over $\F_{p^2}$, which holds through $p=61$ and
	first fails at $p=73$. The obstruction is one of field of definition, not
	of the graph. The successor engine removes it by working, at each vertex,
	over the extension $\F_{p^{2m}}$ with $m$ the least common multiple of the
	Frobenius cycle lengths on the six Weierstrass points, while keeping every
	isomorphism-class key in $\F_{p^2}$ through the absolute Igusa invariants.
	Concretely, at each vertex the engine
	\begin{enumerate}[label=\textup{(\roman*)},leftmargin=2.4em]
		\item computes the Frobenius cycle type on the six Weierstrass points and
		sets $m$ to its least common multiple;
		\item enumerates the fifteen quadratic splittings over $\F_{p^{2m}}$ and
		forms the Richelot codomains;
		\item keys each codomain by its absolute Igusa invariants, which lie in
		$\F_{p^2}$; and
		\item accumulates the out-edge row.
	\end{enumerate}
	Each step is finite-field arithmetic; none consults the automorphic side.
	All certificates of \S\ref{ss:pillars} are unchanged, and two arithmetic
	anchors are imposed at every prime: the fixed-point count of $R(\pi)$ must
	equal the closed formula of \cite[Thm.~2, Rem.~3]{IbukiyamaKatsura}, and
	the trace identity $\Tr R(\pi)=L_0(p)+d(p)$ of
	Theorem~\ref{thm:structural-trace} must hold. The successor engine
	reproduces the matrices of \S\ref{ss:engine} entry by entry at the
	smallest primes before entering new territory.
	
	\subsection{The split-side pipelines}
	\label{ss:split-pipelines}
	
	The factorization certificate is two-sided. On the elliptic side, the
	weight-$2$ eigenvalue systems and their Fricke signs are computed from the
	supersingular $2$-isogeny Brandt matrix together with the involution
	$j\mapsto j^{(p)}$, and the weight-$4$ systems and signs from Manin
	symbols \cite{Manin1972,Cremona1997} for $\Gamma_0(p)$ with Merel's
	Heilbronn matrices \cite{Merel1994}, the sign read off the $U_p$-eigenvalue
	$-\eps_p\,p$. Both pipelines are validated against the closed dimension
	formulas at every prime. The weight-$2$ pipeline reuses the finite-field
	and supersingular-curve utilities of the graph package; the weight-$4$
	pipeline is a separate Manin-symbol implementation.
	
	The characteristic polynomial of each graph is then divided exactly by
	$(x-15)$ times the Saito--Kurokawa and Yoshida products built from the
	orbit polynomials of these data, as in
	Conjecture~\ref{conj:eigenvalue-sign}; the Yoshida blocks are realized as
	characteristic polynomials of Kronecker sums of companion matrices. The
	quotient is required to be an exact square of degree $2\deg\Va_p$ times a
	factor $N_p$ of degree $\delta(p)$.
	
	\subsection{Reduction to a finite certificate}
	\label{ss:finitecert}
	
	The proof of Theorem~\ref{thm:certified-spectra} separates into a
	deductive step, proved here, and a finite list of integer identities
	supplied by the archived data. The deductive step rests on two statements.
	
	\begin{lemma}	\label{lem:onemod}
		Let $M\in M_n(\Z)$ be diagonalizable over $\C$ and let $R\in M_n(\Z)$ be an
		involution with $RM=MR$. Let $f$ be an irreducible factor of $\charp(M)$
		with multiplicity $m_f$, and let $h_f\in\Q[x]$ be the idempotent with
		$h_f\equiv1\pmod{f^{m_f}}$ and $h_f\equiv0\pmod{\charp(M)/f^{m_f}}$. Then
		$h_f(M)$ is the projector onto the $f$-isotypic subspace
		$V_f=\ker f(M)^{m_f}$, an $M$- and $R$-invariant rational subspace of
		dimension $m_f\deg f$; the trace $\Tr\bigl(R\,h_f(M)\bigr)$ is an integer
		of absolute value at most $m_f\deg f\le n$; and for any prime $q>2n$ at
		which the denominators of $h_f$ are invertible, its residue modulo $q$
		determines it.
	\end{lemma}
	
	\begin{proof}
		$V_f$ is rational, $M$- and $R$-invariant, of dimension $m_f\deg f$, and
		$h_f(M)$ is the identity on $V_f$ and zero on the sum of the other
		isotypic subspaces. The restriction $R|_{V_f}$ is an involution of a
		rational space, so its trace is the integer
		$\dim V_f^+-\dim V_f^-$, of absolute value at most $\dim V_f$. An integer
		in $[-n,n]$ is determined by its residue modulo any $q>2n$.
	\end{proof}
	
	\begin{proposition}	\label{prop:reduction}
		Let $n=h_2(p)$. Suppose given $M\in M_n(\Z)$, weights $e\in\Z_{>0}^n$, a
		permutation matrix $R$ with underlying permutation $\sigma$, and the
		weight-$2$ and weight-$4$ eigenvalue data with their Fricke signs. Assume
		\begin{enumerate}[label=\textup{(\alph*)},leftmargin=2.4em]
			\item $\sum_jM_{ij}=15$ and $e_jM_{ij}=e_iM_{ji}$ for all $i,j$;
			\item $R^2=I$, $RM=MR$, and $e_{\sigma(i)}=e_i$;
			\item $\charp(M)=(x-15)\cdot P_{\mathrm{SK}}(x)\cdot
			P_{\mathrm{Y}}(x)\cdot Q(x)^2\cdot N_p(x)$ exactly in $\Z[x]$, where
			$P_{\mathrm{SK}}$ and $P_{\mathrm{Y}}$ are assembled from the given data
			as prescribed by Conjecture~\ref{conj:eigenvalue-sign}, $Q$ is monic, and
			$\deg N_p=\delta(p)$;
			\item for every irreducible $f\mid\charp(M)$,
			$\Tr\bigl(R\,h_f(M)\bigr)=s_f\deg f$, where $h_f$ is as in
			Lemma~\ref{lem:onemod} and $s_f$ is the signed multiplicity the predicted
			eigenvalue--sign multiset assigns to $f$.
		\end{enumerate}
		Then $(M,R)$ is simultaneously diagonalizable over $\R$ and, for every
		eigenvalue $\lambda$ and each choice of sign,
		\[
		\dim\ker(M-\lambda)\cap\ker(R-\eps)
		=\#\{\kappa:\kappa=\lambda,\ \eps_\kappa=\eps\},
		\]
		where the right side counts the predicted multiset. In particular the
		conclusion of Theorem~\ref{thm:certified-spectra} holds at $p$. Summing
		\textup{(d)} over all $f$ gives $\Tr R=\sum_\kappa\eps_\kappa$, the trace
		identity of Theorem~\ref{thm:structural-trace}.
	\end{proposition}
	
	\begin{proof}
		By \textup{(a)} and Lemma~\ref{lem:mestre}(c), $M$ is diagonalizable with
		real spectrum; by \textup{(b)}, $R$ is a diagonalizable involution
		commuting with $M$, so the pair is simultaneously diagonalizable. Write
		$m_\lambda^\pm$ for the $(\lambda,\pm1)$-multiplicities. The
		$\pm1$-eigenspaces of $R$ are rational subspaces, and $M$ restricted to
		each is a rational matrix, so $m_\lambda^\pm=m_{\lambda'}^\pm$ for every
		Galois conjugate: both functions are constant on the roots of each
		irreducible $f$. Hypothesis \textup{(c)} gives
		$m_\lambda^++m_\lambda^-=\#\{\kappa=\lambda\}$ at every $\lambda$. For an
		irreducible $f$, $\Tr(R\,h_f(M))=\sum_{f(\lambda)=0}(m_\lambda^+-m_\lambda^-)
		=\deg f\cdot(m_\lambda^+-m_\lambda^-)$ at any root, so hypothesis
		\textup{(d)} gives $m_\lambda^+-m_\lambda^-=s_f$ there; the predicted
		multiset assigns the same signed count to every root of $f$ block by
		block, by construction, so $s_f$ equals the predicted difference at each
		root. The sum and the difference determine $m_\lambda^\pm$.
	\end{proof}
	
	Hypotheses \textup{(a)}--\textup{(d)} are finitely many equalities between
	integers computable from $M,R,e$ and the block polynomials in exact
	arithmetic; by Lemma~\ref{lem:onemod} the traces in \textup{(d)} may be
	certified modulo a single prime exceeding $2h_2(p)$, provided the
	projector denominators are invertible there. None of the hypotheses
	presupposes any automorphic structure: \textup{(c)} is an exact division
	of integer polynomials and \textup{(d)} a finite list of integer traces,
	both mechanical. The content of the proposition is the passage from these
	checkable identities to the simultaneous eigenspace structure, the
	sign-refined conclusion of Conjecture~\ref{conj:eigenvalue-sign}, which
	no polynomial identity supplies by itself.
	
	\subsection{Worked example: the Ihara pair at $p=19$}
	\label{ss:ihara19}
	
	The principal genus at $p=19$ has $h_2=10$ classes: three products
	(both supersingular $j$-invariants, $j=7$ and $j=18$, lie in $\F_{19}$)
	and seven Jacobians, of which exactly one Galois-conjugate pair
	$\{J_5,J_6\}$ lies over $\F_{19^2}\setminus\F_{19}$; hence $\Tr R(\pi)=8$
	and $T_1(19)=9$. The characteristic polynomial of $B_2(2)$ contains the
	squared factor $(x+2)^2$: this is the block $\Va_{19}$. The eigenvalue
	$-2$ has the two-dimensional eigenspace spanned by
	\[
	w^+=(4,-8,16,10,2,-16,-4,6,3,3),\qquad
	w^-=e_{J_5}-e_{J_6},
	\]
	and $R(\pi)$ acts as the transposition of $J_5$ and $J_6$: it fixes $w^+$
	and negates $w^-$. That $w^-$ is an eigenvector is visible in the matrix
	itself: column $J_5$ minus column $J_6$ of $B_2(2)$ equals
	$-2\,(e_{J_5}-e_{J_6})$. The two eigensystems through $-2$ therefore agree
	at $T(2)$ (and, by Theorem~\ref{thm:certified-spectra}, at every $T(n)$
	with $19\nmid n$) and are separated only by their Atkin--Lehner signs
	$\pm1$: the first type-$\Va$ pair, exhibited on the graph.
	
	\subsection{Worked example: the matrix at $p=61$}
	\label{ss:p61}
	
	The engine assembles the full $128\times128$ matrix $B_2(2)$ at $p=61$
	($15$ product vertices and $113$ Jacobian vertices) and passes the three
	certificates (P1)--(P3) of \S\ref{ss:pillars}. Factoring its
	characteristic polynomial over $\Q$ gives
	\begin{equation}\label{eq:charp61}
		\charp B_2(2)
		=(x-15)\,
		\underbrace{m_{g_1}\!(x{-}6)\,m_{g_2}\!(x{-}6)}_{\text{Saito--Kurokawa},\
			\deg 15}\,
		\underbrace{y_6(x)\,y_{27}(x)}_{\text{Yoshida},\ \deg 33}\,
		\Va_{61}(x)^2\,(x+7),
	\end{equation}
	where $g_1,g_2$ run over the two Galois orbits of
	$S_4(\Gamma_0(61))^{\mathrm{new}}$ (degrees $6$ and $9$) and $m_g$ is the
	minimal polynomial of $a_2(g)$: the Saito--Kurokawa factors are exactly
	$m_g(x-6)$, confirming the shift $+6$ of
	Proposition~\ref{prop:arthur} directly from the elliptic data. The Yoshida
	factors $y_6,y_{27}$ have degrees $6$ and $27$, and $\deg\Va_{61}=39$, so
	the block degrees are $(1,15,33,2\cdot39,1)$, summing to $128$, and the
	trace of $B_2(2)$ is $126$. The fixed-point count $\Tr R(\pi)(61)=38$ matches the
	closed formula of Theorem~\ref{thm:structural-trace}, equivalently
	$T_1(61)=\tfrac12(128+38)=83$.
	
	The general-type factor is the single linear factor $N_{61}(x)=x+7$, and
	$-7$ is a root of none of the other factors, so it is a simple eigenvalue
	isolated by the geometry alone: the engine counts $(2,2)$-isogeny edges
	over $\F_{61^2}$ and never uses the value $-7$. Thus
	\begin{equation}\label{eq:lambda61}
		\lambda^\ast=-7\quad\text{is realized geometrically at }p=61.
	\end{equation}
	This agrees with two further determinations of the same number. First, on
	the paramodular side, $\delta(61)=1$ and $\dim S_3^+(K(61))=0$
	\cite[Cor.~5.3]{Ibukiyama2019quinary} force a single non-lift $f_{61}$
	with Fricke sign $\eps_{61}(f_{61})=-1$, unconditionally; the trace
	identity of Theorem~\ref{thm:structural-trace},
	$\Tr R(\pi)(61)-L_0(61)=38-37=1=-\eps_{61}(f_{61})$, then confirms the
	general-type sign clause of Conjecture~\ref{conj:eigenvalue-sign} at
	$p=61$ (its $R(\pi)$-sign is $+1$), and Proposition~\ref{prop:frame}
	bounds the geometric eigenvalue by $|\lambda^\ast|\le14$. Second, the
	paramodular value $a_2(f_{61})=-7$ of \cite{PoorYuen2015}, obtained by
	Borcherds products, matches the geometric root of \eqref{eq:charp61}. All
	three agree.
	
	\subsection{The full range}
	\label{ss:full-range}
	
	Table~\ref{tab:ext} records the certified block degrees and the
	general-type factor at every prime. The general-type factor $N_p$ is empty
	for $p\le59$; in the certified range it has degree one at
	$p=61,73,79,89,113$, degree two at $p=97,101,103,131,137$, degree three
	at $p=109,127$, and degree four at $p=139,149$. The rational values are
	$-7,-6,-5,-4,-3$ at $p=61,73,79,89,113$; the quadratic factors occur over
	$\Q(\sqrt5)$ and $\Q(\sqrt2)$; and the cubic blocks at $p=109,127$ have
	cyclic Galois group with Hecke fields $\Q(\zeta_7)^+$ and
	$\Q(\zeta_9)^+$.
	
	\begin{table}[ht]
		\centering
		\small
		\begin{tabular}{@{}rrrrrrl@{}}
			\toprule
			$p$ & $h_2$ & $\dim\snew_4$ & $\mathrm{cross}(p)$ & $\deg\Va_p$ &
			$\delta(p)$ & $N_p$\\
			\midrule
			41 & 50 & 10 & 21 & 9 & 0 & ---\\
			43 & 55 & 10 & 16 & 14 & 0 & ---\\
			47 & 72 & 11 & 32 & 14 & 0 & ---\\
			53 & 93 & 13 & 29 & 25 & 0 & ---\\
			59 & 125 & 14 & 50 & 30 & 0 & ---\\
			67 & 166 & 16 & 41 & 54 & 0 & ---\\
			71 & 198 & 17 & 72 & 54 & 0 & ---\\
			73 & 204 & 18 & 46 & 69 & 1 & $x+6$\\
			79 & 256 & 19 & 67 & 84 & 1 & $x+5$\\
			83 & 296 & 20 & 85 & 95 & 0 & ---\\
			89 & 352 & 22 & 92 & 118 & 1 & $x+4$\\
			97 & 436 & 24 & 85 & 162 & 2 & $x^2+9x+19$\\
			101 & 493 & 25 & 121 & 172 & 2 & $x^2+7x+11$\\
			103 & 518 & 25 & 110 & 190 & 2 & $x^2+9x+19$\\
			107 & 581 & 26 & 132 & 211 & 0 & ---\\
			109 & 600 & 27 & 111 & 229 & 3 & $x^3+10x^2+31x+29$\\
			113 & 668 & 28 & 132 & 253 & 1 & $x+3$\\
			127 & 918 & 31 & 165 & 359 & 3 & $x^3+9x^2+24x+19$\\
			131 & 1008 & 32 & 221 & 376 & 2 & $x^2+6x+7$\\
			137 & 1134 & 34 & 193 & 452 & 2 & $x^2+4x+2$\\
			139 & 1179 & 34 & 202 & 469 & 4 & $x^4+14x^3+63x^2+110x+61$\\
			149 & 1433 & 37 & 243 & 574 & 4 & $x^4+6x^3+10x^2+3x-1$\\
			\bottomrule
		\end{tabular}
		\caption{The extension range $41\le p\le149$. The range $11\le p\le37$
			appears in the appendix; $p=61$ is detailed in \S\ref{ss:p61}. At every
			listed prime the exact division and the sign certificate passed, and the
			trace anchors matched \cite[Thm.~2, Rem.~3]{IbukiyamaKatsura} and
			Theorem~\ref{thm:structural-trace}.}
		\label{tab:ext}
	\end{table}
	
	\subsection{Reproducibility and the artifact}
	\label{ss:reproducibility}
	
	The computational data and code accompanying this paper are archived as
	version~3.0.0 of the Richelot--Brandt artifact
	\cite{artifact}. The immutable archive is available at the version DOI
	\doi{10.5281/zenodo.21927983}; the concept DOI
	\doi{10.5281/zenodo.20798967} resolves to the latest version; the development repository is
	\url{https://github.com/dangthanhhung/richelot-brandt}. The archived file
	\texttt{richelot-brandt-artifact-v3.0.0.zip} has SHA-256 digest
	\begin{center}
		\texttt{a68cc8bb01455e6c587cc47757f10fbb\allowbreak 304d0362e32d1f9d895e8f20b597f0a4}
	\end{center}
	The archive
	contains the source code, the per-prime graph data, the weight-$2$ and
	weight-$4$ modular data, the final assembly ledger, and a SHA-256 manifest
	of every file in the archive.
	
	For each prime $11\le p\le149$, the archived graph record contains the
	integer out-edge matrix $M_p$, the automorphism weights, and the
	permutation $\sigma_p$ inducing $R(\pi)$. The archived modular records
	contain the weight-$2$ and weight-$4$ orbit polynomials together with the
	Atkin--Lehner signs used to assemble the Saito--Kurokawa and Yoshida
	blocks. The assembly ledger records the resulting block degrees and the
	general-type factor $N_p$.
	
	The verification separates three layers. First, the graph record is
	checked in exact arithmetic: every row of $M_p$ sums to $15$, the Mestre
	symmetry $e_jM_{ij}=e_iM_{ji}$ holds, the mass identity is satisfied,
	$\sigma_p$ is an involution preserving the weights, and $R(\pi)$ commutes
	with $M_p$. Second, the characteristic polynomial is recomputed and
	divided exactly by the Eisenstein, Saito--Kurokawa, and Yoshida blocks;
	the residual is required to be $\Va_p(x)^2N_p(x)$. Third, for every
	irreducible factor $f$, the signed isotypic trace
	$\operatorname{Tr}(R\,h_f(M_p))$ is checked modulo the prime
	$q=2^{61}-1$, where $h_f$ is the rational projector onto the
	$f$-isotypic subspace. Since $2h_2(p)\le2866<q$ throughout the verified
	range and the required cofactor inversions are recorded, these modular
	traces determine the integer signed traces and hence the
	$R(\pi)$-sign decomposition.
	
	A reference verifier implementing these checks is included as
	\texttt{verify\_all.py}. It reads only the frozen JSON
	records and does not invoke the graph generator. The production checks use
	exact integer and polynomial arithmetic; the external comparison with the
	quinary database of Assaf--Ladd--Rama--Tornar\'ia--Voight is reported
	separately as corroboration and is not an input to the proof of
	Theorem~\ref{thm:certified-spectra}.
	
	\begin{remark}\label{rem:trustedbase}
		The proof of Theorem~\ref{thm:certified-spectra} uses explicit certificates: the archived artifact supplies,
		for each prime, the finite data $M_p,R_p,e$ and the block polynomials, and the hypotheses of Proposition~\ref{prop:reduction} are integer identities replayed from these data in exact
		arithmetic. The computational input consists of the archived certificates with their
		SHA-256 manifest, the verifying code, the published closed formulas invoked in the
		trace comparison and in \S\ref{ss:extension}, and, for the identification of $M_p$ with $B_2(2)$,
		the specification of \S\ref{ss:engine}; the generating implementation is not part of it.
		The identification is over-determined by the two closed-formula anchors of \S\ref{ss:extension},
		the calibrations 	at $p=11$, and the external comparison of \S\ref{ss:external}.
		For 	$11\le p\le37$ the matrices and weights are printed in the appendix,
		so hypothesis \textup{(a)} and the fixed-point counts of Lemma~\ref{lem:fixed-point-trace} are
		checkable from the paper alone; 	hypotheses \textup{(b)}--\textup{(d)} require in addition the
		involutions $\sigma_p$ and the split-side block data, printed in the
		per-prime records of the artifact. The verifying code is small
		and separate from the engines: a reference verifier using only exact
		integer, rational, and polynomial arithmetic replays hypotheses
		\textup{(a)}--\textup{(d)} of Proposition~\ref{prop:reduction} directly
		from the archived data, and shares no code with the generating
		implementation; a FLINT-based routine performs the same replay at the
		largest levels. The computational dependence of the theorem is thereby
		isolated, not eliminated: short of formal verification, hypotheses
		\textup{(a)}--\textup{(d)} must be checked by machine, and
		Theorem~\ref{thm:certified-spectra} is a computer-assisted result in the
		standard sense. The generating algorithm itself is specified in
		\S\ref{ss:engine} and \S\ref{ss:extension}, so the archived matrices can
		be regenerated independently of the artifact.
	\end{remark}
	
	\subsection{External corroboration}
	\label{ss:external}
	
	The general-type data of Theorem~\ref{thm:certified-spectra} admit an
	independent check. The database of
	Assaf--Ladd--Rama--Tornar\'ia--Voight \cite{ALRTV}, computed from
	orthogonal modular forms on quinary lattices, records the Hecke
	eigensystems of the weight-$(3,0)$ paramodular non-lifts of type G at
	every level $N<1000$ under the operator $T_{2,1}$. At every prime
	$11\le p\le149$ the characteristic polynomial of $T_{2,1}$ on the type-G
	newforms of \cite{ALRTV} equals the general-type factor $N_p$ of
	Theorem~\ref{thm:certified-spectra}, including $N_p=1$ at the seventeen
	primes with $\delta(p)=0$; the comparison script accompanies the artifact
	as certificate P4. The graph thus recovers, from $(2,2)$-isogenies alone,
	data obtained in \cite{ALRTV} by an entirely different route; conversely,
	\cite{ALRTV} confirms every general-type block of Table~\ref{tab:ext}. At
	$p=61,73,79$ these blocks also match the paramodular eigenvalue
	computations of \cite{PSY2024}.
	
	This comparison is corroborating evidence, not a proof input: the proof of
	Theorem~\ref{thm:certified-spectra} does not use it. The Atkin--Lehner
	sign convention of \cite{ALRTV} relative to the paramodular Fricke sign
	used here has not been pinned, so the comparison is stated at the level of
	eigenvalue polynomials only.
	
	\subsection{Proof of Theorem~\ref{thm:certified-spectra}}
	\label{ss:proof-certified}
	
	\begin{proof}[Proof of Theorem~\ref{thm:certified-spectra}]
		At every prime the archived data artifact supplies the matrix $M_p$, the
		involution $R_p$, the weights, and the split-side block polynomials with
		their Fricke signs; hypotheses \textup{(a)}--\textup{(d)} of
		Proposition~\ref{prop:reduction} are checked from these data in exact
		arithmetic: \textup{(c)} by exact division over $\Q$, and
		\textup{(d)} through Lemma~\ref{lem:onemod} modulo the Mersenne prime
		$2^{61}-1$, sufficient since $2h_2(p)\le2866<2^{61}-1$.
		Proposition~\ref{prop:reduction} then yields the factorization and the
		signs. For $11\le p\le37$ and $p=61$ the matrices and factorizations are
		those of \S\ref{ss:engine}--\S\ref{ss:p61} and the appendix, replayed
		under the certificates P1--P3; the Ihara separation at $p=19$ is
		\S\ref{ss:ihara19}; summing \textup{(d)} over all factors recovers
		$\Tr R(\pi)$, evaluated against \cite[Thm.~2]{IbukiyamaKatsura} and
		Theorem~\ref{thm:structural-trace}; and the $p=11$ catalogue matches the
		theta computation of \cite[\S4]{Ibukiyama2018conj}. The runs were
		performed in full on the reference machine at every prime, replicated end
		to end on a second machine through $p=103$; \S\ref{ss:external} is the
		independent external comparison.
	\end{proof}
	
	\section{Consequences: type numbers, Ihara pairs, and the spectral profile}
	\label{sec:consequences}
	
	The structural trace identity of Section~\ref{sec:trace} and the certified
	spectra of Section~\ref{sec:certificates} combine to give three direct
	consequences: a paramodular-dimension formula for the principal-genus type
	number, the appearance of Ihara pairs in weight $(3,0)$, and a precise
	description of how the Richelot spectrum departs from its dimension-one
	analogue.
	
	\subsection{The type number in paramodular data}
	\label{ss:typenum-consequence}
	
	We now interpret Corollary~\ref{cor:typenum} from the graph,
	automorphic, and arithmetic viewpoints. Three separately defined quantities meet in this formula: the $\pi$-orbit count on the graph $G_p$, the automorphic trace of
	Theorem~\ref{thm:structural-trace}, and the Ibukiyama--Katsura lattice count of
	$\F_p$-rational superspecial points \cite[Thm.~2, Rem.~3]{IbukiyamaKatsura}. They arise from the combinatorial, automorphic, and arithmetic sides respectively. Their equality is a
	theorem (Theorem~\ref{thm:structural-trace} with \cite[Thm.~2]{IbukiyamaKatsura}), not a coincidence; what makes the agreement a nontrivial check is that the three are computed through separate pipelines (graph enumeration, closed dimension formulas, printed
	lattice tables) and agree at every prime tested (\S\ref{ss:pillars}).
	
	\subsection{Ihara pairs in weight $(3,0)$}
	\label{ss:ihara-consequence}
	
	\begin{corollary}	For each prime $p$ at which a type-$\Va$ parameter occurs, the
		type-$\Va$ representations furnish Ihara pairs: pairs of Hecke
		eigensystems agreeing at every $T(n)$ with $p\nmid n$, separated only
		by $R(\pi)$ into opposite Atkin--Lehner signs. This is unconditional
		(Theorem~\ref{thm:occurrence}). That the number of such pairs equals
		$\deg\Va_p$ is part of Conjecture~\ref{conj:eigenvalue-sign} and
		holds at the certified primes by Theorem~\ref{thm:certified-spectra}.
		The first pair occurs at $p=19$, the least prime with
		$\deg\Va_p\ge1$.
	\end{corollary}
	
	\begin{proof}
		By Theorem~\ref{thm:occurrence}, each occurring type-$\Va$ parameter
		contributes an isotypic space of dimension $2d_p\ge2$ on which every
		$T(n)$ with $p\nmid n$ acts by a scalar and on which the
		$R(\pi)$-eigenspaces for $+1$ and $-1$ each have dimension $d_p$; any
		choice of a $+1$ and a $-1$ eigenline is an Ihara pair. The count
		$\deg\Va_p$ is the degree clause of the conjecture, certified in the
		finite range; the $p=19$ case is displayed explicitly in
		\S\ref{ss:ihara19}.
	\end{proof}
	
	The pairing is the concrete import of the eigenvalue-level refinement. A
	dimension count assigns the two members of a pair a single multiplicity
	and cannot tell them apart; the same is true of every Hecke operator
	$T(n)$ with $p\nmid n$, which acts on both by the same scalar. Only
	$R(\pi)$, the operator outside the away-from-$p$ Hecke algebra,
	distinguishes them. This is the phenomenon Ihara observed in weight
	$(8,8)$ \cite{Ihara1967}, here realized in the geometric weight $(3,0)$ of
	abelian surfaces.
	
	\subsection{The spectral profile}
	\label{ss:spectral-profile}
	
	The spectrum also settles one point that the dimension-one analogy might
	suggest otherwise. The supersingular $\ell$-isogeny graphs of elliptic
	curves are Ramanujan \cite{Pizer1990}, and this expansion underlies their
	use in isogeny-based hashing \cite{CGL2009}. The Richelot graph $G_p$ is
	not: the bound of Proposition~\ref{prop:frame} places its non-Perron
	spectrum in $(-15,15)$, wider than the Ramanujan bound $2\sqrt{14}\approx
	7.48$ for a $15$-regular graph, and the second-largest eigenvalue exceeds
	$2\sqrt{14}$ at every prime in the certified range. The excess is the
	Yoshida and general-type part of the factorization of
	Conjecture~\ref{conj:eigenvalue-sign}, governed by the elliptic Hecke
	eigenvalues at $2$ rather than by any expansion property.
	
	This is consistent with the parallel line of work on spectral gaps:
	Aikawa, Tanaka and Yamauchi \cite{AikawaTanakaYamauchi} bound the
	eigenvalues of the $\ell$-marked superspecial graphs in every dimension
	$g\ge2$ via the Kazhdan constant of the symplectic group, and Katsura and
	Yamauchi \cite{KatsuraYamauchi} refine this for abelian surfaces through
	the Jacquet--Langlands correspondence. Those are bounds on the spectrum of
	a family of graphs; Codogni and Lido \cite{CodogniLido} give a general
	spectral framework for isogeny graphs. The present paper is orthogonal: it factors the
	characteristic polynomial of a single graph exactly, matching each
	eigenvalue to a named automorphic type, and the exact factorization
	records precisely the Yoshida and general-type excess that a spectral-gap
	bound treats in the aggregate.
	
	\begin{remark}\label{rem:nonramanujan}
		The non-Ramanujan conclusion is a statement about the weighted adjacency
		operator $B_2(2)$ with its Brandt normalization, whose Perron eigenvalue
		is $15$ and whose non-Perron eigenvalues are the subject of
		Conjecture~\ref{conj:eigenvalue-sign}. It does not require the conjecture:
		the second-largest eigenvalue exceeds $2\sqrt{14}$ already at the
		certified primes, by direct inspection of the factorizations of
		Section~\ref{sec:certificates}.
	\end{remark}
	
	\section{The Fricke sign bias}
	\label{sec:signbias}
	
	The last invariant of the catalogue is the Fricke sign bias
	\begin{equation}\label{eq:dp-def}
		d(p)=\delta(p)-2\dim S_3^+(K(p))
		=\#\{\text{non-lifts with }\eps_p=-1\}
		-\#\{\text{non-lifts with }\eps_p=+1\},
	\end{equation}
	the second equality holding because every weight-$3$ Gritsenko lift \cite{Gritsenko1995} is
	paramodular-minus \cite{IPY2013}. The main theorems do not depend on this
	section: Theorems~\ref{thm:transfer-framework} and
	\ref{thm:structural-trace} hold unconditionally for all $p\ge11$ and
	$p\ge7$ respectively, and their proofs never use the sign of
	$d(p)$ (Remark~\ref{rem:dp-status}). What $d(p)\ge0$ governs is one
	downstream refinement: the Atkin--Lehner $\pm$-split of the general-type
	block, i.e.\ whether the non-lift eigenforms sit in the Fricke-minus
	eigenspace. Within this section we prove the exact closed formula for $d(p)$ as an
	evaluation of Ibukiyama's non-principal-genus trace (Proposition~\ref{prop:exactmass}), and
	the inequality $d(p)\ge0$ itself for every prime (Theorem~\ref{thm:bias}): the closed
	formula exhibits the $p^{3/2}$ term of $d(p)$ as a positive multiple of
	$B_{2,\chi}$, and the effective lower bound \eqref{eq:B2-lower-bound} then
	makes the sign unconditional.
	
	\subsection{The fixed-point trace and the exact formula for $d(p)$}
	\label{ss:mass}
	
	The involution $R(\pi)$ permutes the vertex set of $G_p$; write $\sigma$
	for this permutation, induced by composing a lattice class with the
	norm-$p$ ideal $\pi$.
	
	\begin{lemma}\label{lem:fixed-point-trace}
		For every prime $p\ge7$,
		\begin{equation}\label{eq:fixtrace}
			\Tr\!\left(R(\pi)\mid M_{0,0}(\Upr(p))\right)
			=2T_1(p)-h_2(p)
			=\sum_{v\in V(G_p)}[\sigma v=v],
		\end{equation}
		where $T_1(p)$ is the principal-genus type number.
	\end{lemma}
	
	\begin{proof}
		The identity $\Tr R(\pi)=2T_1-h_2$ is orbit counting for an involution on
		a finite set, the number of orbits being $\frac12(h_2+\mathrm{Fix})$
		\cite[Thm.~2]{IbukiyamaKatsura}. On the certified graphs the counts are
		$\mathrm{Fix}=5,4,8,8,14,18,18,11$ for $p=11,\dots,37$, and
		$\mathrm{Fix}=38$ at $p=61$.
	\end{proof}
	
	\begin{proposition}	\label{prop:exactmass}
		For every prime $p\ge7$,
		\begin{equation}\label{eq:exactmass}
			d(p)=\Tr R_{\mathrm{npg}}(\pi)-1-\dim S_4^-(p)
			=-\!\!\sum_{v\in N_p}\!\!\eps_p(v),
		\end{equation}
		the last sum running over the general-type non-lift classes $N_p$ with
		$\eps_p(v)$ their paramodular Fricke eigenvalues. With
		$\Tr R_{\mathrm{npg}}(\pi)$ given by
		\cite[Thm.~5.2]{Ibukiyama2019quinary} and
		$\dim S_4^-=\tfrac12\bigl(\dim\snew_4-\tfrac12\nu(p)\bigr)$
		(Lemma~\ref{lem:elliptic-input}), this is the explicit class-number
		identity: for $p\equiv1\pmod4$,
		\[
		d(p)=\frac{9-2s}{96}B_{2,\chi}
		+\frac{5}{16}h(\sqrt{-p})
		+\frac18h(\sqrt{-2p})
		+\frac{3+s}{12}h(\sqrt{-3p})
		-\frac{p+7}{8},
		\]
		and for $p\equiv3\pmod4$,
		\[
		d(p)=\frac{1}{96}B_{2,\chi}
		+\frac{13-5s}{16}h(\sqrt{-p})
		+\frac18h(\sqrt{-2p})
		+\frac1{12}h(\sqrt{-3p})
		-\frac{p+5}{8},
		\]
		with $s=\kro{2}{p}$. There is no multiplicity denominator: each non-lift eigensystem contributes a single line to $S_3(K(p))$ (newform multiplicity one \cite[Prop.~10.1]{RoesnerWeissauer2021}), so the sum is a plain signed count.
	\end{proposition}
	
	\begin{proof}
		The first equality is the non-principal trace evaluation of
		\eqref{eq:npg-trace-d}: $H_{\mathrm{npg}}=1+\dim S_3(K(p))$ and
		$\Tr R_{\mathrm{npg}}=H_{\mathrm{npg}}-2\dim S_3^+(K(p))$, so
		$\Tr R_{\mathrm{npg}}-1-\dim S_4^-
		=\dim S_3(K(p))-\dim S_4^--2\dim S_3^+(K(p))=d(p)$. For the second,
		$W_p$ on $S_3(K(p))$ has trace
		$2\dim S_3^+(K(p))-\dim S_3(K(p))=1-\Tr R_{\mathrm{npg}}$. For the third,
		the lift subspace of dimension $\dim S_4^-$ lies in $S_3^-(K(p))$
		\cite{IPY2013} and contributes $-\dim S_4^-$ to the trace; removing it
		leaves $-\sum_{v\in N_p}\eps_p(v)$. The explicit displays substitute
		\eqref{eq:TOH1}--\eqref{eq:TOH3}
		together with $\dim S_4^-=\tfrac12\bigl(\dim\snew_4-\tfrac12\nu(p)\bigr)$
		and the classical $\dim\snew_4=\tfrac14\bigl(p-2+\kro{-1}{p}\bigr)$; they were verified
		in exact arithmetic against the first equality for all primes
		$7\le p\le2500$, with anchors $d(11)=0$, $d(61)=1$, and, in the regime
		$\dim S_3^+(K(p))=0$, $d(167)=2$, $d(173)=4$, $d(197)=5$, $d(227)=4$,
		$d(233)=9$ \cite{artifact}.
	\end{proof}
	
	Every quantity in \eqref{eq:exactmass} is a theorem: the identity is
	unconditional. By Lemma~\ref{lem:twist}, its class-number terms are
	precisely the CM-class contributions to the non-principal-genus trace, so
	the sign question for $d(p)$ is a domination question among CM classes on
	that genus.
	
	\subsection{The optimal-embedding dictionary}
	\label{ss:twist}
	
	We record the dictionary ``torsion class $\leftrightarrow$ CM discriminant
	$\leftrightarrow$ class number $h(-D)$'' that interprets the individual
	terms of the trace formulas, with the discriminant corrections the literal
	embedding counts require.
	
	\begin{lemma}	\label{lem:twist}
		Let $B=B_{p,\infty}$ and let $\Onorm$ be a maximal order. For an imaginary
		quadratic order $R$ of discriminant $-D$, the number of optimal embeddings
		$R\hookrightarrow\Onorm$ up to $\Onorm^\times$-conjugacy is
		\[
		m(-D)=h(-D)\cdot\prod_{q\mid p\infty}
		\Bigl(1-\Bigl\{\frac{-D}{q}\Bigr\}\Bigr),
		\qquad
		\Bigl\{\frac{-D}{p}\Bigr\}=
		\begin{cases}\kro{-D}{p} & p\nmid D,\\ 0 & p\mid D,\end{cases}
		\]
		where $\kro{\cdot}{p}$ is the Legendre symbol and $h(-D)$ the class number
		of $R$ (with $h(-D)=0$ unless $-D\equiv0,1\bmod4$). The torsion classes
		correspond to the following data:
		\begin{itemize}
			\item[$(\ell{=}2)$] order-$4$ torsion $\leftrightarrow\Q(\sqrt{-1})$,
			contributing through $-D=-4p$ when $p\equiv1\bmod4$ and through $-D=-p$
			when $p\equiv3\bmod4$; the $h(\sqrt{-2p})$ term is the
			$\Q(\sqrt{-2})$-embedding, gated by the local symbol $\kro{-2}{p}$;
			\item[$(\ell{=}3)$] order-$3$ (and $6$) torsion $\leftrightarrow\Q(\sqrt{-3})$,
			contributing through $-D=-3p$, gated by $\kro{-3}{p}$.
		\end{itemize}
	\end{lemma}
	
	\begin{proof}
		The formula is Eichler's optimal-embedding count for definite quaternion
		orders \cite[Ch.~III]{Vigneras}, with the local factor at the ramified
		prime $p$ given by the Eichler symbol; at the infinite ramified place the
		factor is $1$ for imaginary $R$. The discriminant clauses are the
		condition $-D\equiv0,1\bmod4$: when $p\equiv1\bmod4$, $-p\equiv3\bmod4$
		is not a discriminant and one uses the fundamental discriminant $-4p$,
		with the conductor relation of \cite[Cor.~7.28]{Cox}; this is the content
		of \eqref{eq:nu-principal}. The order-$3$ versus order-$6$ distinction is
		the unit group of $\Z[\omega]$ and does not affect the embedding count.
	\end{proof}
	
	\subsection{The Fricke-minus regime}
	\label{ss:uncond-regime}
	
	\begin{corollary}\label{thm:uncond-regime}
		$\dim S_3^+(K(p))=0$ for every prime $p<167$
		\cite[Cor.~5.3]{Ibukiyama2019quinary},
		\cite[Prop.~8.2]{Ibukiyama-dim}; hence every weight-$3$ non-lift of level
		$p<167$ is Fricke-minus and $d(p)=\delta(p)\ge0$ unconditionally there,
		both vanishing for $p\le59$. Beyond that range, $\dim S_3^+(K(p))=1$
		exactly for $p\in\{167,173,197,223,233,239,251,271,277,281,313,331,337\}$
		and $=2$ exactly for
		$p\in\{227,257,263,269,283,349,379,409,421\}$
		\cite[Prop.~8.2]{Ibukiyama-dim}, and direct evaluation of
		\eqref{eq:exactmass} gives $d(p)\ge0$ at every one of these primes as
		well. In particular $d(p)=\delta(p)$ at $p=61,73,79$: the Calabi--Yau
		eigenforms lie in the Fricke-minus eigenspace.
	\end{corollary}
	
	\begin{proof}
		The vanishing and the two lists are Ibukiyama's theorems, via
		$\dim S_3^+(K(p))=\tfrac12\bigl(H_{\mathrm{npg}}(p)-\Tr
		R_{\mathrm{npg}}(\pi)\bigr)$; the rest is \eqref{eq:dp-def} together with
		$\delta(p)\ge0$ (Lemma~\ref{lem:positivity}) and the tabulated values of
		Proposition~\ref{prop:exactmass}.
	\end{proof}
	
	For $p\ge167$ some non-lifts are Fricke-plus, and $d(p)\ge0$ becomes the
	statement that the Fricke-minus classes still dominate, which the
	following theorem establishes at every prime.
	
	\subsection{The Bernoulli term and the remaining sign question}
	\label{ss:bernoulli-term}
	
	Let $\chi$ be the primitive real quadratic character of conductor $D$ and
	let $K=\mathbf Q(\sqrt D)$.  With the convention
	\[
	L(1-n,\chi)=-\frac{B_{n,\chi}}{n}\qquad(n\geq1),
	\]
	the factorization $\zeta_K(s)=\zeta(s)L(s,\chi)$ at $s=-1$ gives the
	exact identity
	\begin{equation}\label{eq:B2-zetaK}
		\zeta_K(-1)=\frac{B_{2,\chi}}{24}.
	\end{equation}
	The completed functional equation for the even primitive character $\chi$ is
	\[
	\left(\frac{D}{\pi}\right)^{s/2}\Gamma\!\left(\frac{s}{2}\right)L(s,\chi)
	=
	\left(\frac{D}{\pi}\right)^{(1-s)/2}
	\Gamma\!\left(\frac{1-s}{2}\right)L(1-s,\chi).
	\]
	Evaluating it at $s=-1$ and using
	$L(-1,\chi)=-B_{2,\chi}/2$ yields
	\begin{equation}\label{eq:B2-functional-equation}
		B_{2,\chi}=\frac{D^{3/2}}{\pi^2}L(2,\chi).
	\end{equation}
	In particular $B_{2,\chi}>0$.  Notice that \emph{$L(2,\chi)$ occurs in
		the numerator}, not its inverse.
	
	For reference, the Euler factors give the uniform lower bound
	\begin{equation}\label{eq:L2-lower-bound}
		L(2,\chi)
		=\prod_{\ell\nmid D}(1-\chi(\ell)\ell^{-2})^{-1}
		\geq\prod_{\ell}(1+\ell^{-2})^{-1}
		=\frac{\zeta(4)}{\zeta(2)}=\frac{\pi^2}{15}.
	\end{equation}
	Consequently
	\begin{equation}\label{eq:B2-lower-bound}
		B_{2,\chi}\geq\frac{D^{3/2}}{15}.
	\end{equation}
	Here a prime dividing $D$ contributes the Euler factor $1$, which is also at
	least $(1+\ell^{-2})^{-1}$; thus no extra correction at $2$ is needed.
	
	Equations~\eqref{eq:B2-zetaK}--\eqref{eq:B2-lower-bound} identify the
	positive leading contribution in the exact formula of
	Proposition~\ref{prop:exactmass}.  Every class-number coefficient
	there is positive in both residue classes, so the only negative term is
	linear in $p$, and \eqref{eq:B2-lower-bound} converts the sign question
	into an explicit inequality.
	
	\begin{theorem}[Positivity of the bias]\label{thm:bias}
		For every prime $p$, $d(p)\ge0$: among the weight-$3$ general-type
		paramodular non-lifts of level $p$, those with Fricke eigenvalue $-1$
		are at least as numerous as those with eigenvalue $+1$; equivalently
		$\dim S_3^+(K(p))\le\tfrac12\delta(p)$. Moreover $d(p)\gg p^{3/2}$.
	\end{theorem}
	
	\begin{proof}
		For $p\le5$, $\dim S_3^{+}(K(p))=0$ by \cite[Prop.~8.2]{Ibukiyama-dim},
		and the Gritsenko-lift argument of Lemma~\ref{lem:positivity} applies at
		every prime, so $d(p)=\delta(p)\ge0$. Assume $p\ge7$.
		Dropping the class-number terms of Proposition~\ref{prop:exactmass},
		whose coefficients are positive, and applying \eqref{eq:B2-lower-bound}
		with $D=p$, respectively $D=4p$, leaves
		\[
		d(p)\ge\frac{9-2s}{96}\cdot\frac{p^{3/2}}{15}-\frac{p+7}{8}
		\quad(p\equiv1\bmod4),
		\qquad
		d(p)\ge\frac{1}{96}\cdot\frac{8p^{3/2}}{15}-\frac{p+5}{8}
		\quad(p\equiv3\bmod4),
		\]
		and both right-hand sides are positive for every prime $p>673$. For
		$p\le673$ the exact evaluation of \eqref{eq:exactmass} gives
		$d(p)\ge0$ directly; the verification recorded in the proof of
		Proposition~\ref{prop:exactmass} covers all primes $7\le p\le2500$
		\cite{artifact}. The growth statement follows from the same lower
		bound.
	\end{proof}
	
	\begin{remark}\label{rem:sign-accessible}
		The defect $d(p)$ is the difference of two quantities of order $p^2$
		with the same leading constant, and magnitude bounds of the kind
		proving Lemma~\ref{lem:positivity} cannot fix the sign of the
		$O(p^{3/2})$ residue. Proposition~\ref{prop:exactmass} changes the
		frame: it exhibits the residue itself as a positive multiple of
		$B_{2,\chi}$ plus positive class-number terms minus a linear term,
		and the effective bound \eqref{eq:B2-lower-bound}, an explicit form
		of the positivity of $\zeta_K(-1)$, then decides the sign. The
		obstruction dissolved with the formula, not with a new estimate.
		Ibukiyama's sign theorem \cite[Thm.~7.1]{Ibukiyama-dim} gives
		$(-1)^k\bigl(\dim S_k^{+}(K(p))-\dim S_k^{-}(K(p))\bigr)\ge0$ for every
		$k\ge3$ and every prime; at $k=3$, since
		$\dim S_3^{-}(K(p))-\dim S_3^{+}(K(p))=\dim S_4^{-}+d(p)$, it amounts to
		$d(p)\ge-\dim S_4^{-}$, and its proof runs through the same
		non-principal trace and its $B_{2,\chi}$ main term.
		Theorem~\ref{thm:bias} removes the lift cushion: the Fricke-minus
		classes dominate already among the non-lifts. The
		trace identity, the transfer theorem, and the certified factorization
		theorem do not use Theorem~\ref{thm:bias}.
	\end{remark}
	
	\section{Concluding remarks}
	\label{sec:concl}
	

\begin{thebibliography}{99}
\bibitem[ALRTV]{ALRTV} E. Assaf, W. Ladd, G. Rama, G. Tornar\'ia, J. Voight, \emph{A database of paramodular forms from quinary orthogonal modular forms}, in: LuCaNT: LMFDB, Computation, and Number Theory, Contemp. Math. \textbf{796}, Amer. Math. Soc., 2024, \doi{10.1090/conm/796/16004}; arXiv:2308.09824. Data: \url{https://github.com/assaferan/omf5_data}.
\bibitem[AikawaTanakaYamauchi]{AikawaTanakaYamauchi} Y. Aikawa, R. Tanaka, T. Yamauchi, \emph{Isogeny graphs on superspecial abelian varieties: eigenvalues and connection to Bruhat--Tits buildings}, Canad. J. Math. \textbf{76} (2024), no.~6, 1891--1916.
\bibitem[AtkinLehner]{AtkinLehner} A.~O.~L. Atkin, J. Lehner, \emph{Hecke operators on $\Gamma_0(m)$}, Math. Ann. \textbf{185} (1970), 134--160.
\bibitem[AtkinMorain1993]{AtkinMorain1993} A.~O.~L. Atkin, F. Morain, \emph{Elliptic curves and primality proving}, Math. Comp. \textbf{61} (1993), no.~203, 29--68.
\bibitem[CGL2009]{CGL2009} D.~X. Charles, K.~E. Lauter, E.~Z. Goren, \emph{Cryptographic hash functions from expander graphs}, J. Cryptology \textbf{22} (2009), no.~1, 93--113.
\bibitem[ChanGan2015]{ChanGan2015} P.-S. Chan, W.~T. Gan, \emph{The local Langlands conjecture for $\GSp(4)$ III: stability and twisted endoscopy}, J. Number Theory \textbf{146} (2015), 69--133.
\bibitem[CodogniLido]{CodogniLido} G. Codogni, G. Lido, \emph{Spectral theory of isogeny graphs}, Commun. Number Theory Phys. \textbf{17} (2023), no.~1, 173--248.
\bibitem[Cox]{Cox} D.~A. Cox, \emph{Primes of the Form $x^2+ny^2$}, 2nd ed., Wiley, 2013.
\bibitem[Cremona1997]{Cremona1997} J.~E. Cremona, \emph{Algorithms for modular elliptic curves}, 2nd ed., Cambridge Univ. Press, Cambridge, 1997.
\bibitem[DPRT]{DPRT} N. Dummigan, A. Pacetti, G. Rama, G. Tornar\'ia, \emph{Quinary forms and paramodular forms}, Math. Comp. \textbf{93} (2024), no.~348, 1805--1858.
\bibitem[Dang2025]{Dang2025} H. T. Dang, D. V. Nguyen, \emph{Exact Stratification and Affine Mass Formulas for Split Richelot Data over Finite Fields}, arXiv:2607.03376 (2026).
\bibitem[FloritSmith]{FloritSmith} E. Florit, B. Smith, \emph{Automorphisms and isogeny graphs of abelian varieties, with applications to the superspecial Richelot isogeny graph}, in: Arithmetic, Geometry, Cryptography, and Coding Theory 2021, Contemp. Math. \textbf{779}, Amer. Math. Soc., 2022, pp.~103--132.
\bibitem[GanTantono]{GanTantono} W.~T. Gan, W. Tantono, \emph{The local Langlands conjecture for $\mathrm{GSp}(4)$, II: The case of inner forms}, Amer. J. Math. \textbf{136} (2014), 761--805.
\bibitem[GeeTaibi]{GeeTaibi} T. Gee, O. Ta\"ibi, \emph{Arthur's multiplicity formula for $\mathrm{GSp}_4$ and restriction to $\mathrm{Sp}_4$}, J. \'Ec. polytech. Math. \textbf{6} (2019), 469--535.
\bibitem[Gritsenko1995]{Gritsenko1995} V. Gritsenko, \emph{Arithmetical lifting and its applications}, in: Number Theory (Paris 1992--93), London Math. Soc. Lecture Note Ser. \textbf{215}, Cambridge Univ. Press, 1995, pp.~103--126.
\bibitem[Hashimoto1980]{Hashimoto1980} K. Hashimoto, \emph{On Brandt matrices associated with the positive definite quaternion Hermitian forms}, J. Fac. Sci. Univ. Tokyo \textbf{27} (1980), 227--245.
\bibitem[HashimotoIbukiyama]{HashimotoIbukiyama} K. Hashimoto, T. Ibukiyama, \emph{On class numbers of positive definite binary quaternion Hermitian forms}, J. Fac. Sci. Univ. Tokyo \textbf{27} (1980), 549--601.
\bibitem[IKO1986]{IKO1986} T. Ibukiyama, T. Katsura, F. Oort, \emph{Supersingular curves of genus two and class numbers}, Compositio Math. \textbf{57} (1986), 127--152.
\bibitem[IPY2013]{IPY2013} T. Ibukiyama, C. Poor, D. S. Yuen, \emph{Jacobi forms that characterize paramodular forms}, Abh. Math. Semin. Univ. Hambg. \textbf{83} (2013), no.~1, 111--128.
\bibitem[Ibukiyama-dim]{Ibukiyama-dim} T. Ibukiyama, \emph{Dimensions of paramodular forms and compact twist modular forms with involutions}, J. Math. Soc. Japan, advance publication (2025), \doi{10.2969/jmsj/93479347}, arXiv:2208.13578v2.
\bibitem[Ibukiyama2007compact]{Ibukiyama2007compact} T. Ibukiyama, \emph{Paramodular forms and compact twist}, in: Automorphic Forms on $GSp(4)$, Proc.\ of the 9th Autumn Workshop on Number Theory, Hakuba, 2006, pp.~37--48.
\bibitem[Ibukiyama2007dim]{Ibukiyama2007dim} T. Ibukiyama, \emph{Dimension formulas of Siegel modular forms of weight $3$ and supersingular abelian surfaces}, Proc. of the 4th Spring Conf. on Modular Forms (2007), 39--60.
\bibitem[Ibukiyama2018conj]{Ibukiyama2018conj} T. Ibukiyama, \emph{Conjectures on correspondence of symplectic modular forms of middle parahoric type and Ihara lifts}, Res. Math. Sci. \textbf{5} (2018), no.~2, Paper No.~18, 36 pp.
\bibitem[Ibukiyama2018type]{Ibukiyama2018type} T. Ibukiyama, \emph{Type numbers of quaternion hermitian forms and supersingular abelian varieties}, Osaka J. Math. \textbf{55} (2018), 369--384.
\bibitem[Ibukiyama2019quinary]{Ibukiyama2019quinary} T. Ibukiyama, \emph{Quinary lattices and binary quaternion hermitian lattices}, Tohoku Math. J. \textbf{71} (2019), 207--220.
\bibitem[IbukiyamaKatsura]{IbukiyamaKatsura} T. Ibukiyama, T. Katsura, \emph{On the field of definition of superspecial polarized abelian varieties and type numbers}, Compositio Math. \textbf{91} (1994), 37--46.
\bibitem[Igusa1960]{Igusa1960} J.-I. Igusa, \emph{Arithmetic variety of moduli for genus two}, Ann. of Math. (2) \textbf{72} (1960), 612--649.
\bibitem[Ihara1967]{Ihara1967} Y. Ihara, \emph{Hecke polynomials as congruence $\zeta$ functions in the elliptic modular case}, Ann. of Math. \textbf{85} (1967), 267--295.
\bibitem[JordanZaytman]{JordanZaytman} B. W. Jordan, Y. Zaytman, \emph{Isogeny graphs of superspecial abelian varieties and Brandt matrices}, J. Math. Soc. Japan, advance publication (2026), \doi{10.2969/jmsj/91629162}; arXiv:2005.09031.
\bibitem[KatsuraTakashima]{KatsuraTakashima} T. Katsura, K. Takashima, \emph{Counting Richelot isogenies between superspecial abelian surfaces}, in: ANTS XIV -- Proc.\ of the Fourteenth Algorithmic Number Theory Symposium, Open Book Ser. \textbf{4}, no.~1, Math. Sci. Publ., 2020, pp.~283--300.
\bibitem[KatsuraYamauchi]{KatsuraYamauchi} T. Katsura, T. Yamauchi, \emph{Isogeny graphs associated to Moret--Bailly families of supersingular abelian surfaces}, (2026), arXiv:2606.21859.
\bibitem[Manin1972]{Manin1972} Yu.~I. Manin, \emph{Parabolic points and zeta-functions of modular curves}, Izv. Akad. Nauk SSSR Ser. Mat. \textbf{36} (1972), 19--66.
\bibitem[Merel1994]{Merel1994} L. Merel, \emph{Universal Fourier expansions of modular forms}, in: On Artin's conjecture for odd $2$-dimensional representations, Lecture Notes in Math. \textbf{1585}, Springer, 1994, pp.~59--94.
\bibitem[Ogg1974]{Ogg1974} A. P. Ogg, \emph{Hyperelliptic modular curves}, Bull. Soc. Math. France \textbf{102} (1974), 449--462.
\bibitem[PSY2024]{PSY2024} C. Poor, J. Shurman, D. S. Yuen, \emph{Eigenvalues and congruences for the weight $3$ paramodular nonlifts of levels $61$, $73$, and $79$}, J. Korean Math. Soc. \textbf{61} (2024), no.~5, 997--1033.
\bibitem[Pizer1990]{Pizer1990} A.~K. Pizer, \emph{Ramanujan graphs and Hecke operators}, Bull. Amer. Math. Soc. (N.S.) \textbf{23} (1990), no.~1, 127--137.
\bibitem[PoorYuen2015]{PoorYuen2015} C. Poor, D. S. Yuen, \emph{Paramodular cusp forms}, Math. Comp. \textbf{84} (2015), no.~293, 1401--1438.
\bibitem[Pratt1975]{Pratt1975} V.~R. Pratt, \emph{Every prime has a succinct certificate}, SIAM J. Comput. \textbf{4} (1975), no.~3, 214--220.
\bibitem[RSST1997]{RSST1997} N. Robertson, D.~P. Sanders, P. Seymour, R. Thomas, \emph{The four-colour theorem}, J. Combin. Theory Ser.~B \textbf{70} (1997), no.~1, 2--44.
\bibitem[RobertsSchmidt]{RobertsSchmidt} B. Roberts, R. Schmidt, \emph{Local Newforms for $GSp(4)$}, Lecture Notes in Math. 1918, Springer, 2007.
\bibitem[RobertsSchmidtTables]{RobertsSchmidtTables} B. Roberts, R. Schmidt, \emph{Tables for representations of $\mathrm{GSp}(4)$}, companion tables to \cite{RobertsSchmidt}.
\bibitem[Roesner2016]{Roesner2016} M. R\"osner, \emph{Parahoric restriction for $GSp(4)$ and the inner cohomology of Siegel modular threefolds}, PhD thesis, Heidelberg, 2016.
\bibitem[RoesnerWeissauer2021]{RoesnerWeissauer2021} M. R\"osner, R. Weissauer, \emph{Global liftings between inner forms of $\GSp(4)$}, J. Number Theory \textbf{263} (2024), 79--138; arXiv:2103.14715.
\bibitem[Schmidt2018]{Schmidt2018} R. Schmidt, \emph{Packet structure and paramodular forms}, Trans. Amer. Math. Soc. \textbf{370} (2018), 3085--3112.
\bibitem[Shimura1963]{Shimura1963} G. Shimura, \emph{Arithmetic of alternating forms and quaternion hermitian forms}, J. Math. Soc. Japan \textbf{15} (1963), 33--65.
\bibitem[Vigneras]{Vigneras} M.-F. Vign\'eras, \emph{Arithm\'etique des Alg\`ebres de Quaternions}, Lecture Notes in Math. 800, Springer, 1980.
\bibitem[Yamauchi1973]{Yamauchi1973} M. Yamauchi, \emph{On the traces of Hecke operators for a normalizer of $\Gamma_0(N)$}, J. Math. Kyoto Univ. \textbf{13} (1973), 403--411.
\bibitem[artifact]{artifact} H. T. Dang, \emph{Computational artifact for ``The Richelot isogeny graph as a Brandt matrix''} (scripts, matrices, and certificates), version 3.0.0, Zenodo, 2026, \doi{10.5281/zenodo.21927983}.
\bibitem[vanHoften]{vanHoften} P. van Hoften, \emph{A geometric Jacquet--Langlands correspondence for paramodular Siegel threefolds}, Math. Z. \textbf{299} (2021), 2029--2061.
\end{thebibliography}
\end{document}
